% Options for packages loaded elsewhere
\PassOptionsToPackage{unicode}{hyperref}
\PassOptionsToPackage{hyphens}{url}
\PassOptionsToPackage{dvipsnames,svgnames,x11names}{xcolor}
%
\documentclass[
  11pt,
  letterpaper,
]{article}
\usepackage{amsmath,amssymb}
\usepackage{iftex}
\ifPDFTeX
  \usepackage[T1]{fontenc}
  \usepackage[utf8]{inputenc}
  \usepackage{textcomp} % provide euro and other symbols
\else % if luatex or xetex
  \usepackage{unicode-math} % this also loads fontspec
  \defaultfontfeatures{Scale=MatchLowercase}
  \defaultfontfeatures[\rmfamily]{Ligatures=TeX,Scale=1}
\fi
\usepackage{lmodern}
\ifPDFTeX\else
  % xetex/luatex font selection
\fi
% Use upquote if available, for straight quotes in verbatim environments
\IfFileExists{upquote.sty}{\usepackage{upquote}}{}
\IfFileExists{microtype.sty}{% use microtype if available
  \usepackage[]{microtype}
  \UseMicrotypeSet[protrusion]{basicmath} % disable protrusion for tt fonts
}{}
\makeatletter
\@ifundefined{KOMAClassName}{% if non-KOMA class
  \IfFileExists{parskip.sty}{%
    \usepackage{parskip}
  }{% else
    \setlength{\parindent}{0pt}
    \setlength{\parskip}{6pt plus 2pt minus 1pt}}
}{% if KOMA class
  \KOMAoptions{parskip=half}}
\makeatother
\usepackage{xcolor}
\usepackage[margin=0.76in]{geometry}
\usepackage{color}
\usepackage{fancyvrb}
\newcommand{\VerbBar}{|}
\newcommand{\VERB}{\Verb[commandchars=\\\{\}]}
\DefineVerbatimEnvironment{Highlighting}{Verbatim}{commandchars=\\\{\}}
% Add ',fontsize=\small' for more characters per line
\newenvironment{Shaded}{}{}
\newcommand{\AlertTok}[1]{\textcolor[rgb]{1.00,0.00,0.00}{\textbf{#1}}}
\newcommand{\AnnotationTok}[1]{\textcolor[rgb]{0.38,0.63,0.69}{\textbf{\textit{#1}}}}
\newcommand{\AttributeTok}[1]{\textcolor[rgb]{0.49,0.56,0.16}{#1}}
\newcommand{\BaseNTok}[1]{\textcolor[rgb]{0.25,0.63,0.44}{#1}}
\newcommand{\BuiltInTok}[1]{\textcolor[rgb]{0.00,0.50,0.00}{#1}}
\newcommand{\CharTok}[1]{\textcolor[rgb]{0.25,0.44,0.63}{#1}}
\newcommand{\CommentTok}[1]{\textcolor[rgb]{0.38,0.63,0.69}{\textit{#1}}}
\newcommand{\CommentVarTok}[1]{\textcolor[rgb]{0.38,0.63,0.69}{\textbf{\textit{#1}}}}
\newcommand{\ConstantTok}[1]{\textcolor[rgb]{0.53,0.00,0.00}{#1}}
\newcommand{\ControlFlowTok}[1]{\textcolor[rgb]{0.00,0.44,0.13}{\textbf{#1}}}
\newcommand{\DataTypeTok}[1]{\textcolor[rgb]{0.56,0.13,0.00}{#1}}
\newcommand{\DecValTok}[1]{\textcolor[rgb]{0.25,0.63,0.44}{#1}}
\newcommand{\DocumentationTok}[1]{\textcolor[rgb]{0.73,0.13,0.13}{\textit{#1}}}
\newcommand{\ErrorTok}[1]{\textcolor[rgb]{1.00,0.00,0.00}{\textbf{#1}}}
\newcommand{\ExtensionTok}[1]{#1}
\newcommand{\FloatTok}[1]{\textcolor[rgb]{0.25,0.63,0.44}{#1}}
\newcommand{\FunctionTok}[1]{\textcolor[rgb]{0.02,0.16,0.49}{#1}}
\newcommand{\ImportTok}[1]{\textcolor[rgb]{0.00,0.50,0.00}{\textbf{#1}}}
\newcommand{\InformationTok}[1]{\textcolor[rgb]{0.38,0.63,0.69}{\textbf{\textit{#1}}}}
\newcommand{\KeywordTok}[1]{\textcolor[rgb]{0.00,0.44,0.13}{\textbf{#1}}}
\newcommand{\NormalTok}[1]{#1}
\newcommand{\OperatorTok}[1]{\textcolor[rgb]{0.40,0.40,0.40}{#1}}
\newcommand{\OtherTok}[1]{\textcolor[rgb]{0.00,0.44,0.13}{#1}}
\newcommand{\PreprocessorTok}[1]{\textcolor[rgb]{0.74,0.48,0.00}{#1}}
\newcommand{\RegionMarkerTok}[1]{#1}
\newcommand{\SpecialCharTok}[1]{\textcolor[rgb]{0.25,0.44,0.63}{#1}}
\newcommand{\SpecialStringTok}[1]{\textcolor[rgb]{0.73,0.40,0.53}{#1}}
\newcommand{\StringTok}[1]{\textcolor[rgb]{0.25,0.44,0.63}{#1}}
\newcommand{\VariableTok}[1]{\textcolor[rgb]{0.10,0.09,0.49}{#1}}
\newcommand{\VerbatimStringTok}[1]{\textcolor[rgb]{0.25,0.44,0.63}{#1}}
\newcommand{\WarningTok}[1]{\textcolor[rgb]{0.38,0.63,0.69}{\textbf{\textit{#1}}}}
\usepackage{longtable,booktabs,array}
\usepackage{calc} % for calculating minipage widths
% Correct order of tables after \paragraph or \subparagraph
\usepackage{etoolbox}
\makeatletter
\patchcmd\longtable{\par}{\if@noskipsec\mbox{}\fi\par}{}{}
\makeatother
% Allow footnotes in longtable head/foot
\IfFileExists{footnotehyper.sty}{\usepackage{footnotehyper}}{\usepackage{footnote}}
\makesavenoteenv{longtable}
\usepackage{graphicx}
\makeatletter
\def\maxwidth{\ifdim\Gin@nat@width>\linewidth\linewidth\else\Gin@nat@width\fi}
\def\maxheight{\ifdim\Gin@nat@height>\textheight\textheight\else\Gin@nat@height\fi}
\makeatother
% Scale images if necessary, so that they will not overflow the page
% margins by default, and it is still possible to overwrite the defaults
% using explicit options in \includegraphics[width, height, ...]{}
\setkeys{Gin}{width=\maxwidth,height=\maxheight,keepaspectratio}
% Set default figure placement to htbp
\makeatletter
\def\fps@figure{htbp}
\makeatother
\setlength{\emergencystretch}{3em} % prevent overfull lines
\providecommand{\tightlist}{%
  \setlength{\itemsep}{0pt}\setlength{\parskip}{0pt}}
\setcounter{secnumdepth}{5}
\ifLuaTeX
\usepackage[bidi=basic]{babel}
\else
\usepackage[bidi=default]{babel}
\fi
\babelprovide[main,import]{american}
% get rid of language-specific shorthands (see #6817):
\let\LanguageShortHands\languageshorthands
\def\languageshorthands#1{}
\usepackage{amsmath,amssymb,mathtools,booktabs,longtable,array,graphicx,microtype,float,xcolor,tcolorbox}
\definecolor{ink}{HTML}{172033}
\definecolor{accent}{HTML}{345B83}
\definecolor{soft}{HTML}{EEF4F8}
\definecolor{warn}{HTML}{FAF3E8}
\definecolor{open}{HTML}{F3EFF8}
\newcommand{\BKM}{\mathrm{BKM}}
\newcommand{\E}{\mathcal E}
\newcommand{\dd}{\mathrm d}
\newcommand{\sech}{\operatorname{sech}}
\newcommand{\avg}[1]{\left\langle #1\right\rangle}
\newcommand{\vperp}{\varrho_\perp}
\newcommand{\Ruble}{\mathfrak R_c}
\newcommand{\Xscale}{\mathcal X_\sigma}
\newcommand{\Sbar}{\mathcal S}
\newcommand{\rc}{r_c}
\newtcolorbox{thesisbox}{colback=soft,colframe=accent,boxrule=0.8pt,arc=1.2mm,left=3mm,right=3mm,top=2mm,bottom=2mm}
\newtcolorbox{warningbox}{colback=warn,colframe=orange!55!black,boxrule=0.7pt,arc=1.2mm,left=3mm,right=3mm,top=2mm,bottom=2mm}
\newtcolorbox{openbox}{colback=open,colframe=purple!45!black,boxrule=0.7pt,arc=1.2mm,left=3mm,right=3mm,top=2mm,bottom=2mm}
\ifLuaTeX
  \usepackage{selnolig}  % disable illegal ligatures
\fi
\usepackage{bookmark}
\IfFileExists{xurl.sty}{\usepackage{xurl}}{} % add URL line breaks if available
\urlstyle{same}
\hypersetup{
  pdftitle={Causal Scale Dynamics},
  pdfauthor={Thomas Ruble},
  pdflang={en-US},
  colorlinks=true,
  linkcolor={blue},
  filecolor={Maroon},
  citecolor={Blue},
  urlcolor={blue},
  pdfcreator={LaTeX via pandoc}}

\title{Causal Scale Dynamics}
\usepackage{etoolbox}
\makeatletter
\providecommand{\subtitle}[1]{% add subtitle to \maketitle
  \apptocmd{\@title}{\par {\large #1 \par}}{}{}
}
\makeatother
\subtitle{Modular Information Geometry, Horizon Capacity, and the Ruble
Equations}
\author{Thomas Ruble}
\date{21 August 2026}

\begin{document}
\maketitle

{
\hypersetup{linkcolor=}
\setcounter{tocdepth}{3}
\tableofcontents
}
\begin{quote}
\textbf{Working master research note v7.0. Not peer reviewed.} This note
distinguishes \textbf{{[}STANDARD{]}}, \textbf{{[}THEOREM{]}},
\textbf{{[}DEDUCTION{]}}, \textbf{{[}PRINCIPLE{]}},
\textbf{{[}CONDITIONAL{]}}, \textbf{{[}NEGATIVE{]}}, and
\textbf{{[}OPEN{]}} statements. The homogeneous theory is closed by one
explicitly stated physical equivalence principle, not advertised as a
theorem of mathematics alone. The covariant perturbation theory and the
global vacuum sector remain separate research problems. The declared
cosmological baseline is four-dimensional spatially flat FLRW with
radiation, pressureless matter, one collective scale-capacity response,
and an exactly zero residual late-time floor. Current observational
comparisons are background-level or response-level analyses, not an
official joint Boltzmann likelihood.
\end{quote}

\section*{Abstract}\label{abstract}
\addcontentsline{toc}{section}{Abstract}

This note develops a self-contained formulation of gravity and cosmic
history in which causal order, metric scale, modular state geometry,
horizon entropy, and late-time acceleration are treated as different
registers of one structure rather than as unrelated inputs.

Under the standard causality and regularity hypotheses, causal order
determines a Lorentzian spacetime only up to conformal scale. The metric
is represented by a conformal class \([g]\) and a positive scale section

\[
\sigma\in\Gamma(\mathcal E[1]),
\qquad
g_\sigma=\sigma^{-2}\boldsymbol g.
\]

Its two-jet is the scale tractor

\[
I_A=\frac14D_A\sigma.
\]

In four dimensions the trace-free Einstein equation is equivalent to the
scale-transport equation

\[
\boxed{
\left(\nabla_a\nabla_b\sigma+P_{ab}\sigma\right)_0
=\frac{4\pi G}{c^4}\,\sigma T_{ab}^{\circ},
}
\]

while the trace and one scalar lift determine

\[
\boxed{
I^2=\frac{2\pi G}{3c^4}T-\frac{\Lambda_g}{3}.
}
\]

Trace-free stress is therefore the obstruction to parallel transport of
scale; the cosmological constant belongs to a distinct global
calibration channel.

A causal region and faithful state carry another geometry: modular flow,
relative entropy, and its Bogoliubov--Kubo--Mori Hessian. The corrected
theory separates vertical modular automorphism flow from horizontal
deformation of the state family. At a homogeneous codimension-two cut,
the normal plane supplies a canonical chirality operator

\[
Q=P_+-P_-,
\qquad
Q^2=1,
\qquad
JQJ=-Q.
\]

The associated binary exponential family is

\[
\omega_\theta=\frac{e^{\theta Q}}{2\cosh\theta},
\]

with dual information coordinates

\[
\eta=\left\langle Q\right\rangle=\tanh\theta,
\qquad
G^{\mathrm{BKM}}_{\theta\theta}=\frac{\mathrm d\eta}{\mathrm d\theta}=\operatorname{sech}^2\theta.
\]

The state coordinate is soldered to logarithmic Weyl scale

\[
N=\ln\frac{a}{a_c}=-\ln\frac{\sigma}{\sigma_c}
\]

through a measurable Connes-cocycle character. Under the rank-one ratio
and measurability hypotheses, cocycle composition forces

\[
\theta=\varrho_\perp(N-N_c).
\]

The value of \(\varrho_\perp\) is not fixed by Cauchy's equation. The
present theory selects the fundamental normal representation,
\(\varrho_\perp=1\), as a physical structural law rather than as a
continuously fitted dark-energy parameter.

The central closure is the \textbf{Scale--Capacity Equivalence
Principle}. Let \(S_c\) be the Bekenstein--Hawking entropy of the
self-dual causal wall, \(T_c=\hbar c/(2\pi k_BR_c)\) its canonically
normalized horizontal modular temperature, \(V_c\) its areal volume, and
\(G^{\perp}_{NN}\) the BKM norm of the fundamental horizontal scale
tangent. The principle states

\[
\boxed{
\frac{k_B}{S_c}\,G^{\perp}_{NN}(N)
=\operatorname{sech}^2(N-N_c).
}
\]

Equivalently, the entropy-normalized BKM speed at self-duality is unity:

\[
\boxed{
\mathfrak R_c
:=\frac{k_B}{S_c}G^{\perp}_{NN}(N_c)=1.
}
\]

The physical source is the quadratic modular free-energy curvature per
causal-wall volume,

\[
\boxed{
\rho_X(N)=\frac{k_BT_c}{2V_c}G^{\perp}_{NN}(N).
}
\]

In \(3+1\) dimensional flat FLRW, horizon thermodynamics gives

\[
k_BT_c\frac{S_c}{k_B}
=E_{\rm MS,c}
=\rho_{\rm crit,c}V_c.
\]

Consequently,

\[
\boxed{
\rho_X(N)
=\frac12\rho_{\rm crit,c}\operatorname{sech}^2(N-N_c).
}
\]

At the self-dual crossing,

\[
\boxed{
\Omega_{X,c}=\frac12,
\qquad
\rho_X(N_c)=\rho_{\rm ordinary}(N_c).
}
\]

The former free normalization \(r_c=1\) is therefore replaced by the
scale-capacity principle. Relative to dust alone there is a calculable
radiation correction of order \(4\times10^{-4}\) for the benchmark
cosmology.

Separate conservation yields

\[
\boxed{
w_X(N)=-1+\frac23\tanh(N-N_c),
}
\]

and the parameter-free differential invariant

\[
\boxed{
9(1+w_X)^2+6\frac{\mathrm dw_X}{\mathrm dN}=4.
}
\]

The sign change of cosmic acceleration is not inserted as an independent
event. It follows from the competition between ordinary dilution and a
single self-dual susceptibility pulse. The model predicts one observed
acceleration entry, one future exit, and asymptotic coasting when the
global residual floor is exactly zero.

For \(\Omega_{m0}=0.310598\) and \(\Omega_{r0}=9.15\times10^{-5}\), the
capacity-normalized branch predicts

\[
N_c=-0.294007,
\qquad
z_c=0.341793,
\]

\[
q_0=-0.336902,
\qquad
j_0=-0.111246,
\]

\[
z_{\rm acc}=0.785694,
\qquad
\frac{a_{\rm exit}}{a_0}=11.7865.
\]

The previously reported background comparison gives
\(\chi^2\simeq1398.3\) for the rigid history and \(1401.6\) for flat
\(\Lambda\)CDM with the same number of background shape parameters. This
is viability, not discovery. The strongest empirical test is the
redshift-by-redshift invariant above, followed by the prediction that
self-duality, \(w_X=-1\), and ordinary--dark equality coincide.

The theory is closed at the homogeneous level by one equivalence law and
one global sector choice. It does not yet provide a complete covariant
perturbation action. The binary information geometry canonically
generates a reflectionless Witten/Darboux pair, but the embedding of
that pair into scalar, vector, and tensor spacetime perturbations
remains open. That is the next research layer rather than an unfixed
background parameter.

\section{The Ruble Equations}\label{the-ruble-equations}

The theory is most compactly stated as a sequence of definitions and one
physical equivalence principle. The equations are organized by
mathematical type rather than by discovery history.

\subsection{Kinematic and geometric
register}\label{kinematic-and-geometric-register}

\textbf{R1 --- causal scale.}

\[
\boxed{
g_{\rm phys}=\sigma^{-2}\boldsymbol g,
\qquad
N=-\ln\frac{\sigma}{\sigma_c}=\ln\frac{a}{a_c}.
}
\]

The conformal metric \(\boldsymbol g\) fixes causal cones; the scale
section \(\sigma\in\Gamma(\mathcal E[1])\) fixes physical calibration.
\(N\) is logarithmic Weyl scale, not Newtonian absolute time.

\textbf{R2 --- scale tractor.}

\[
\boxed{
I_A=\frac14D_A\sigma.
}
\]

The scale tractor packages \(\sigma\), its first derivative, and the
trace combination of its second derivative.

\textbf{R3 --- tractor transport and norm.}

\[
\boxed{
\left(\nabla_a\nabla_b\sigma+P_{ab}\sigma\right)_0
=\frac{4\pi G}{c^4}\sigma T_{ab}^{\circ},
}
\]

\[
\boxed{
I^2=\frac{2\pi G}{3c^4}T-\frac{\Lambda_g}{3}.
}
\]

The first equation is local and trace-free. The second fixes the scalar
calibration channel up to a global constant.

\subsection{Horizontal state register}\label{horizontal-state-register}

\textbf{R4 --- normal chirality.}

\[
\boxed{
Q=P_+-P_-,
\qquad
Q^2=1,
\qquad
JQJ=-Q.
}
\]

\(P_\pm\) project onto the two null lines of the Lorentzian normal plane
of a codimension-two cut. The homogeneous horizontal response is assumed
to factor through this fundamental chirality quotient.

\textbf{R5 --- binary exponential family.}

\[
\boxed{
\omega_\theta=\frac{e^{\theta Q}}{2\cosh\theta},
\qquad
\Psi(\theta)=\ln(2\cosh\theta).
}
\]

Its dual coordinates and metric are

\[
\boxed{
\eta=\Psi'(\theta)=\tanh\theta,
\qquad
G^{\mathrm{BKM}}_{\theta\theta}=\Psi''(\theta)=\operatorname{sech}^2\theta.
}
\]

The identity

\[
\boxed{
\eta^2+G^{\mathrm{BKM}}_{\theta\theta}=1
}
\]

is the normalized binary moment relation
\(\left\langle Q^2\right\rangle=1\).

\textbf{R6 --- horizontal soldering.}

Under the rank-one ratio and measurability hypotheses, the Connes
cocycle chain rule gives

\[
\theta=\varrho_\perp(N-N_c).
\]

The theory selects the fundamental character

\[
\boxed{\varrho_\perp=1.}
\]

This is a physical representation choice: the fundamental null-normal
character is identified with the fundamental scale/inverse-scale
character. It is not a theorem of Cauchy's equation, and it is not
obtained from an integer restriction on conformal weights.

\subsection{Scale-capacity closure}\label{scale-capacity-closure}

\textbf{R7 --- Ruble scale-capacity number.}

Let \(G^{\perp}_{NN}\) be the physical BKM norm of the selected
horizontal scale tangent and \(S_c\) the horizon entropy of the
self-dual wall. Define

\[
\boxed{
\mathfrak R_c:=\frac{k_B}{S_c}G^{\perp}_{NN}(N_c).
}
\]

The Scale--Capacity Equivalence Principle is

\[
\boxed{\mathfrak R_c=1.}
\]

In canonical binary coordinates this is equivalently

\[
\boxed{
\frac{k_B}{S_c}G^{\perp}_{NN}(N)=\operatorname{sech}^2(N-N_c).
}
\]

It states that the fundamental Weyl translation is isometric, at
self-duality, to the fundamental horizontal state translation when the
BKM metric is normalized by horizon entropy.

\textbf{R8 --- modular free-energy source.}

\[
\boxed{
\rho_X(N)=\frac{k_BT_c}{2V_c}G^{\perp}_{NN}(N).
}
\]

This is the constitutive law. It converts dimensionless relative-entropy
curvature into energy density by the physical modular temperature and
causal-wall volume. The factor \(1/2\) is the Taylor coefficient of
relative entropy at coincidence.

\textbf{R9 --- Hawking--Friedmann conversion.}

For a flat four-dimensional FLRW apparent horizon,

\[
\boxed{
R_c=\frac{c}{H_c},
\qquad
\frac{S_c}{k_B}=\frac{\pi R_c^2c^3}{G\hbar},
\qquad
k_BT_c=\frac{\hbar c}{2\pi R_c},
}
\]

and

\[
\boxed{
k_BT_c\frac{S_c}{k_B}
=E_{\rm MS,c}
=\rho_{\rm crit,c}V_c.
}
\]

\textbf{R10 --- closed source.}

Combining R7--R9 gives

\[
\boxed{
\rho_X(N)=\frac12\rho_{\rm crit,c}\operatorname{sech}^2(N-N_c).
}
\]

At the crossing,

\[
\boxed{
\Omega_{X,c}=\frac12,
\qquad
\rho_X(N_c)=\rho_{\rm ordinary}(N_c).
}
\]

The exact equality is with the complete non-dark sector. Relative to
dust alone,

\[
\frac{\rho_X(N_c)}{\rho_m(N_c)}
=1+\frac{\rho_r(N_c)}{\rho_m(N_c)}+\cdots.
\]

\subsection{Dynamical consequences}\label{dynamical-consequences}

\textbf{R11 --- equation of state.}

Separate conservation,

\[
\rho_X'=-3(1+w_X)\rho_X,
\]

gives

\[
\boxed{
w_X(N)=-1+\frac23\tanh(N-N_c).
}
\]

\textbf{R12 --- shape invariant.}

\[
\boxed{
9(1+w_X)^2+6w_X'=4.
}
\]

Equivalent forms are

\[
\boxed{
\frac{\rho_X}{\rho_*}
+\frac14\left(\frac{\mathrm d\ln\rho_X}{\mathrm dN}\right)^2=1,
}
\]

and

\[
X'=\frac23-\frac32X^2,
\qquad X:=1+w_X.
\]

\textbf{R13 --- self-dual comparison.}

The symmetrized relative entropy of the state and its modular reflection
is

\[
\boxed{
\mathfrak S_J(\theta)=4\theta\tanh\theta.
}
\]

It has a unique global minimum at \(\theta=0\). Along the rigid history,

\[
\boxed{
\mathfrak S_J=6(N-N_c)(1+w_X).
}
\]

Thus relative-entropy positivity fixes the orientation of the \(w_X=-1\)
crossing relative to increasing Weyl scale.

\textbf{R14 --- central blindness.}

\[
\boxed{
\operatorname{Var}(K+\alpha\mathbf1)=\operatorname{Var}(K).
}
\]

The local response is blind to the absolute energy zero. A global scalar
lift remains in \(\Lambda_g\).

\begin{figure}
\centering
\includegraphics[width=0.94\textwidth,height=\textheight]{figures/dependency_v7.pdf}
\caption{Dependency structure of the theory. The scale-capacity
principle is the single new equivalence law closing the homogeneous
amplitude.}
\end{figure}

\subsection{What is law, what is solution data, and what is
measured}\label{what-is-law-what-is-solution-data-and-what-is-measured}

\begin{longtable}[]{@{}
  >{\raggedright\arraybackslash}p{(\columnwidth - 4\tabcolsep) * \real{0.3333}}
  >{\raggedright\arraybackslash}p{(\columnwidth - 4\tabcolsep) * \real{0.3333}}
  >{\raggedright\arraybackslash}p{(\columnwidth - 4\tabcolsep) * \real{0.3333}}@{}}
\toprule\noalign{}
\begin{minipage}[b]{\linewidth}\raggedright
Quantity
\end{minipage} & \begin{minipage}[b]{\linewidth}\raggedright
Type
\end{minipage} & \begin{minipage}[b]{\linewidth}\raggedright
Role
\end{minipage} \\
\midrule\noalign{}
\endhead
\bottomrule\noalign{}
\endlastfoot
\([g]\) & causal/conformal geometry & fixes null cones and causal
order \\
\(\sigma\) & scale section & metric calibration \\
\(N\) & Weyl coordinate & logarithmic scale displacement \\
\(Q\) & normalized normal chirality & fundamental horizontal binary
score \\
\(\theta\) & horizontal state coordinate & relative modular
polarization \\
\(\mathfrak R_c\) & dimensionless equivalence number &
entropy-normalized peak BKM speed; postulated value \(1\) \\
\(N_c\) & solution origin & location of the intrinsic self-dual event
relative to today \\
\(\Omega_{m0},\Omega_{r0}\) & measured state data & determine the
location of the crossing in the closed background \\
\(\Lambda_g\) & global lift/sector & not part of local trace-free
transport; zero in the open-future branch \\
\end{longtable}

The theory therefore contains no continuously fitted dark-history
function and no independent dark amplitude. It does contain one explicit
physical equivalence principle and one global sector choice.

\section{Retyping the primitive
notions}\label{retyping-the-primitive-notions}

\subsection{Causal order is not Newtonian clock
time}\label{causal-order-is-not-newtonian-clock-time}

Several quantities called ``time'' are categorically different:

\begin{longtable}[]{@{}
  >{\raggedright\arraybackslash}p{(\columnwidth - 4\tabcolsep) * \real{0.3333}}
  >{\raggedright\arraybackslash}p{(\columnwidth - 4\tabcolsep) * \real{0.3333}}
  >{\raggedright\arraybackslash}p{(\columnwidth - 4\tabcolsep) * \real{0.3333}}@{}}
\toprule\noalign{}
\begin{minipage}[b]{\linewidth}\raggedright
Symbol/concept
\end{minipage} & \begin{minipage}[b]{\linewidth}\raggedright
Mathematical type
\end{minipage} & \begin{minipage}[b]{\linewidth}\raggedright
Reversal property
\end{minipage} \\
\midrule\noalign{}
\endhead
\bottomrule\noalign{}
\endlastfoot
causal order \(x\prec y\) & partial order on events & orientation is
part of the relation \\
Newtonian/coordinate \(t\) & parameter in equations & often reversed by
\(t\mapsto-t\) \\
proper time \(\tau\) & metric length along a timelike curve & depends on
calibrated metric \\
Weyl scale time \(N=\ln a\) & additive coordinate on positive scale
ratios & reverses between expansion and contraction \\
modular parameter \(s\) & vertical automorphism parameter of a fixed
state & a reversible one-parameter group \\
horizon rapidity \(\eta_A\) & geometric boost coordinate when modular
flow is geometric & vertical clock at a cut \\
horizontal polarization \(\theta\) & coordinate comparing different
states/fibers & state deformation, not elapsed modular time \\
\end{longtable}

The v5.0 notation conflated \(\eta_A\) and \(\theta\). The corrected
theory distinguishes

\[
\boxed{
\frac{\mathrm d\eta_A}{\mathrm dN}=\mu_A(N),
\qquad
\frac{\mathrm d\theta}{\mathrm dN}=\varrho_\perp.
}
\]

The first is a running horizon-clock rate; the second is the proposed
horizontal scale-to-state conversion.

\subsection{Metric geometry is conformal structure plus
scale}\label{metric-geometry-is-conformal-structure-plus-scale}

A conformal manifold is \((M,[g])\), with

\[
g_{ab}\sim\Omega^2g_{ab}.
\]

A positive scale section

\[
\sigma\in\Gamma(\mathcal E[1])
\]

selects the physical metric

\[
g_\sigma=\sigma^{-2}\boldsymbol g.
\]

The pair \(([g],\sigma)\) contains the same information as a metric. The
physical proposal is not that one has discovered an algebraic
decomposition. It is that causal geometry and metric calibration may
obey different dynamical laws and should not be fused before asking what
sources each register.

\subsection{Matter and dark response are not new signed
charges}\label{matter-and-dark-response-are-not-new-signed-charges}

The primitive local source is not a scalar ``gravitational charge.'' In
the conformal description it is the trace-free stress obstructing
parallel scale transport. In the proposed state description it is a
positive information-geometric response. Signed modular orientation
lives in \(\left\langle Q\right\rangle=\tanh\theta\); the susceptibility
is even:

\[
\operatorname{Var}(Q)=\operatorname{sech}^2\theta=\operatorname{sech}^2(-\theta).
\]

This permits conjugate state orientations without introducing positive
and negative gravitational charges.

\section{Conformal gravity as transport and
calibration}\label{conformal-gravity-as-transport-and-calibration}

\subsection{Scale tractor}\label{scale-tractor}

For four-dimensional conformal geometry define

\[
P_{ab}=\frac12\left(R_{ab}-\frac16Rg_{ab}\right),
\qquad
J=P^a{}_a=\frac{R}{6}.
\]

In a chosen representative metric, a standard tractor is written

\[
V^A\simeq(\alpha,\mu_a,\beta),
\]

with tractor norm

\[
h(V,V)=2\alpha\beta+g^{ab}\mu_a\mu_b.
\]

The scale tractor is

\[
\boxed{
I_A=\frac14D_A\sigma
\simeq
\left(
\sigma,
\nabla_a\sigma,
-\frac14(\Delta\sigma+J\sigma)
\right).
}
\]

The almost-Einstein operator is

\[
\boxed{
\mathcal E_{ab}(\sigma)
:=\left(\nabla_a\nabla_b\sigma+P_{ab}\sigma\right)_0.
}
\]

\textbf{{[}STANDARD{]}} On the open set where \(\sigma\neq0\),

\[
\nabla_a^TI_B=0
\Longleftrightarrow
\mathcal E_{ab}(\sigma)=0
\Longleftrightarrow
g_\sigma\text{ is Einstein}.
\]

In the physical scale,

\[
\boxed{I^2=-\frac{R[g_\sigma]}{12}.}
\]

\subsection{Transport and norm
equations}\label{transport-and-norm-equations}

Taking the trace-free part of Einstein's equation gives

\[
G^{\circ}_{ab}=\frac{8\pi G}{c^4}T^{\circ}_{ab}.
\]

The conformal transformation law yields

\[
G^{\circ}_{ab}[g_\sigma]
=2\sigma^{-1}\mathcal E_{ab}(\sigma),
\]

hence

\[
\boxed{
\mathcal E_{ab}(\sigma)
=\frac{4\pi G}{c^4}\sigma T^{\circ}_{ab}.
}
\]

Taking the trace gives

\[
\boxed{
I^2=\frac{2\pi G}{3c^4}T-\frac{\Lambda_{\rm g}}3.
}
\]

The accurate slogan is therefore:

\[
\boxed{
\begin{aligned}
\text{vacuum Einstein geometry}
&=\text{parallel scale transport},\\
\text{noncentral matter stress}
&=\text{transport defect},\\
\text{stress trace}
&=\text{norm variation},\\
\Lambda_{\rm g}
&=\text{global scalar lift}.
\end{aligned}}
\]

A measured \(w_X\neq-1\) is not simply evidence for a new force. In this
language it is evidence that the dark sector contributes to the failure
of scale-tractor parallelism.

\subsection{Why the local equation is blind to the vacuum
zero}\label{why-the-local-equation-is-blind-to-the-vacuum-zero}

Under

\[
T_{ab}\mapsto T_{ab}+\lambda g_{ab},
\]

one has

\[
T^{\circ}_{ab}\mapsto T^{\circ}_{ab}.
\]

Likewise a normalized Gibbs state is invariant under

\[
H\mapsto H+C\mathbf 1.
\]

Thus an additive vacuum offset is central in both languages:

\[
\boxed{
\begin{aligned}
H&\sim H+C\mathbf1,\\
T_{ab}&\sim T_{ab}+\lambda g_{ab},\\
\text{local causal response}&\sim\text{equivalence class modulo the central direction}.
\end{aligned}}
\]

This explains local vacuum blindness. It does not by itself make the
remaining scalar lift radiatively stable. Section 11 returns to the
global completion.

\section{Modular state geometry, self-duality, and the horizontal scale
line}\label{modular-state-geometry-self-duality-and-the-horizontal-scale-line}

\subsection{Vertical modular flow and horizontal state
deformation}\label{vertical-modular-flow-and-horizontal-state-deformation}

For a local von Neumann algebra \(\mathcal A\) and faithful state
\(\omega\), Tomita--Takesaki theory supplies ({[}6{]}) a modular
operator \(\Delta_\omega\), conjugation \(J_\omega\), and automorphism
group

\[
\sigma_s^\omega(A)=\Delta_\omega^{is}A\Delta_\omega^{-is}.
\]

This is \textbf{vertical} motion: an automorphism at fixed state. A
family

\[
N\longmapsto\omega_N
\]

is \textbf{horizontal} motion through state space. A Connes cocycle or
modular Berry connection is required to compare the modular frames of
different fibers.

Schematically,

\[
\boxed{
\partial_NK_N
=[\mathcal A_N,K_N]
+(\mathcal D_NK_N)_\perp
+c_N'\mathbf1.
}
\]

\begin{longtable}[]{@{}
  >{\raggedright\arraybackslash}p{(\columnwidth - 2\tabcolsep) * \real{0.5000}}
  >{\raggedright\arraybackslash}p{(\columnwidth - 2\tabcolsep) * \real{0.5000}}@{}}
\toprule\noalign{}
\begin{minipage}[b]{\linewidth}\raggedright
term
\end{minipage} & \begin{minipage}[b]{\linewidth}\raggedright
mathematical role
\end{minipage} \\
\midrule\noalign{}
\endhead
\bottomrule\noalign{}
\endlastfoot
\([\mathcal A_N,K_N]\) & change of modular frame / vertical gauge \\
\((\mathcal D_NK_N)_\perp\) & physical horizontal noncentral
deformation \\
\(c_N'\mathbf1\) & central normalization shift \\
\end{longtable}

The gravitational boost-charge principle acts on the vertical generator.
Scale capacity acts on the horizontal tangent. The vacuum energy zero is
central.

This distinction is essential. The vertical horizon rapidity satisfies

\[
\frac{\mathrm d\eta_A}{\mathrm dN}=\mu_A=\frac{1-q}{2},
\]

whereas the horizontal state coordinate obeys

\[
\frac{\mathrm d\theta}{\mathrm dN}=\varrho_\perp.
\]

They are not the same parameter.

\subsection{The normal pair and the fundamental
score}\label{the-normal-pair-and-the-fundamental-score}

A spacelike codimension-two cut has a Lorentzian normal plane

\[
N(\Sigma)=L_+\oplus L_-.
\]

Let \(P_\pm\) be the two null-line projectors. The normal chirality
operator

\[
Q=P_+-P_-
\]

satisfies

\[
Q^2=P_++P_-=1.
\]

Normal reflection exchanges the rays and gives

\[
JQJ=-Q.
\]

The full modular generator of a type-III algebra need not have two
eigenvalues. The structural identification is narrower:

\begin{quote}
\textbf{Fundamental normal reduction.} The homogeneous \(J\)-odd
horizontal response is the infrared projection of the full state
deformation onto the normal chirality quotient.
\end{quote}

This is the one physical representation choice underlying the binary
model. It is not a claim that the entire quantum field theory has a
two-level Hilbert space.

\subsection{Cocycle composition and affine
soldering}\label{cocycle-composition-and-affine-soldering}

Assume the reduced relative cocycle depends on two scale sections only
through their ratio \(r=\sigma_2/\sigma_1\) and has one noncentral
generator \(Q\). Connes' chain rule then reduces to

\[
\theta(r_1r_2)=\theta(r_1)+\theta(r_2).
\]

Measurability is sufficient to solve the Cauchy equation:

\[
\boxed{
\theta(r)=-\varrho_\perp\ln r
=\varrho_\perp(N-N_c).
}
\]

This derives the logarithmic form but not the coefficient. The Claude
closure's proposed integrality argument is not valid in ordinary
conformal geometry: density bundles \(\mathcal E[w]\) exist for every
real \(w\), so a real character \(r^{\varrho_\perp}\) is mathematically
admissible. The present theory selects the unit fundamental character as
a physical law,

\[
\boxed{\varrho_\perp=1,}
\]

rather than pretending that Cauchy's equation or conformal
representation theory alone excludes all other real values.

\subsection{Binary information
geometry}\label{binary-information-geometry}

The maximum-entropy family generated by \(Q\) is

\[
\omega_\theta=\frac{e^{\theta Q}}{Z(\theta)},
\qquad
Z(\theta)=2\cosh\theta.
\]

The log-partition potential

\[
\Psi(\theta)=\ln Z(\theta)
\]

generates the mixture coordinate and BKM metric:

\[
\eta=\Psi'(\theta)=\left\langle Q\right\rangle=\tanh\theta,
\]

\[
\boxed{
G^{\mathrm{BKM}}_{\theta\theta}=\Psi''(\theta)=\operatorname{Var}_\theta(Q)=\operatorname{sech}^2\theta.
}
\]

The identity

\[
\eta^2+G^{\mathrm{BKM}}_{\theta\theta}=1
\]

is simply \(\left\langle Q^2\right\rangle=1\).

The full Fisher distance of the crossover is

\[
\boxed{
L_F=\int_{-\infty}^{\infty}\operatorname{sech}\theta\,\mathrm d\theta=\pi.
}
\]

The state traverses the complete binary simplex from one extremal ray to
the other. The width in Weyl scale is fixed by the soldering
coefficient.

\begin{figure}
\centering
\includegraphics[width=0.9\textwidth,height=\textheight]{figures/fisher_complete_traversal.pdf}
\caption{The complete binary traversal. The total Fisher distance is
fixed; the Weyl soldering determines how that distance is distributed
over cosmic scale.}
\end{figure}

\subsection{Modular self-duality}\label{modular-self-duality}

Because \(JQJ=-Q\),

\[
J\omega_\theta J=\omega_{-\theta}.
\]

The symmetrized relative entropy is ({[}19{]})

\[
\boxed{
\mathfrak S_J(\theta)
=S(\omega_\theta\|\omega_{-\theta})
+S(\omega_{-\theta}\|\omega_\theta)
=4\theta\tanh\theta.
}
\]

It is nonnegative, vanishes only at \(\theta=0\), and has

\[
\mathfrak S_J'(0)=0,
\qquad
\mathfrak S_J''(0)=8.
\]

The crossing is therefore intrinsically defined as the unique self-dual
state. Its numerical date relative to today remains solution data
determined by the matter/radiation state and flatness.

The relation

\[
\mathfrak S_J=6(N-N_c)(1+w_X)
\]

later implies that the phantom-to-quintessence orientation is fixed
relative to increasing Weyl scale. It does not derive why the universe
chooses the expanding rather than time-reversed branch.

\subsection{What information geometry
fixes}\label{what-information-geometry-fixes}

Čencov's theorem fixes the classical Fisher metric up to a global scale
under the relevant statistical morphisms. Petz classifies quantum
monotone metrics; monotonicity alone does not uniquely select BKM. The
relative-entropy Hessian selects the BKM metric, and the binary
normalization \(Q^2=1\) fixes the dimensionless shape.

None of those theorems converts dimensionless BKM curvature into energy
density. That conversion is supplied by the scale-capacity equivalence
principle and horizon thermodynamics, not by a hidden normalization
theorem.

\section{The Scale--Capacity Equivalence
Principle}\label{the-scalecapacity-equivalence-principle}

The previous versions contained two numbers:

\[
\varrho_\perp=\frac{\mathrm d\theta}{\mathrm dN}
\]

for the horizontal conversion rate and

\[
\chi_\perp
\]

for the conversion from BKM susceptibility to energy density. Background
data partly degenerate them. The present closure replaces the arbitrary
stiffness by an invariant relation among the state metric, horizon
entropy, modular temperature, and causal volume.

\subsection{Capacity and the physical BKM
norm}\label{capacity-and-the-physical-bkm-norm}

The normalized binary metric

\[
\operatorname{sech}^2\theta
\]

fixes shape but is not the extensive capacity of the whole causal wall.
Write

\[
G^{\perp}_{NN}(N)
=C_{\perp,c}\varrho_\perp^2\operatorname{sech}^2[\varrho_\perp(N-N_c)],
\]

where

\[
C_{\perp,c}=G^{\perp}_{\theta\theta}(0)
\]

is the physical BKM capacity carried by the selected horizontal mode.

The entanglement capacity of a state is

\[
C_E=\operatorname{Var}(K).
\]

It is not equal to entropy for arbitrary states. In controlled
holographic CFTs with an ordinary Einstein-gravity dual and the
preferred gravitational regularization, spherical regions satisfy
({[}26{]})

\[
C_E=S_{EE}/k_B.
\]

Higher-derivative gravity and generic quantum states can change the
ratio. The equality therefore characterizes an Einstein-capacity
universality class rather than a tautology.

\subsection{The invariant Ruble
number}\label{the-invariant-ruble-number}

Define

\[
\boxed{
\mathfrak R_c
:=\frac{k_B}{S_c}G^{\perp}_{NN}(N_c)
=\frac{k_BC_{\perp,c}}{S_c}\varrho_\perp^2.
}
\]

This combination is invariant under rescaling the coordinate \(\theta\):
the BKM metric and coordinate slope transform oppositely. It is the
entropy-normalized squared speed of the horizontal state under one Weyl
e-fold.

The fundamental law is

\[
\boxed{\mathfrak R_c=1.}
\]

Equivalently,

\[
\boxed{
\left.
\Phi^*\left(\frac{k_B}{S_c}G^{\perp}_{\mathrm{BKM}}\right)
\right|_{N_c}
=\mathrm dN^2.
}
\]

This is the stopping principle of the homogeneous theory. It is
analogous in logical status to an equivalence principle or a universal
conversion law: it is not derived from a more primitive metric
convention, and it is directly falsifiable.

Under the fundamental normal representation, \(\varrho_\perp=1\), the
same law implies

\[
C_{\perp,c}=S_c/k_B.
\]

Conversely, if the selected wall mode saturates the Einstein capacity,
the unit Ruble number fixes \(\varrho_\perp=1\).

\subsection{Relative entropy as modular free-energy
curvature}\label{relative-entropy-as-modular-free-energy-curvature}

Let \(\omega_c\) be the self-dual KMS reference state and
\(\mathcal H_c=k_BT_cK_c\) its physical modular Hamiltonian.
Nonequilibrium free energy obeys the exact identity

\[
F_c(\rho)-F_c(\omega_c)
=k_BT_cS(\rho\|\omega_c).
\]

For a neighboring scale state,

\[
S(\omega_{c+\delta N}\|\omega_c)
=\frac12G^{\perp}_{NN}(N_c)\delta N^2+O(\delta N^3).
\]

Therefore the quadratic free-energy curvature is

\[
\lim_{\delta N\to0}
\frac{F_c(\omega_{c+\delta N})-F_c(\omega_c)}{\delta N^2}
=\frac{k_BT_c}{2}G^{\perp}_{NN}(N_c).
\]

The homogeneous source law is defined by distributing that collective
modular free energy over the causal-wall volume:

\[
\boxed{
\rho_X(N)=\frac{k_BT_c}{2V_c}G^{\perp}_{NN}(N).
}
\]

This is not a canonical local scalar kinetic term. \(\theta\) is a
collective state coordinate, and the energy is the free-energy curvature
of the state family. The single-field ghost no-go therefore does not
invalidate the homogeneous construction.

\subsection{Hawking--Friedmann
conversion}\label{hawkingfriedmann-conversion}

For a flat FLRW apparent horizon in four spacetime dimensions,

\[
R_c=\frac{c}{H_c},
\qquad
A_c=4\pi R_c^2,
\qquad
V_c=\frac{4\pi}{3}R_c^3.
\]

Use the Bekenstein--Hawking entropy

\[
\frac{S_c}{k_B}=\frac{A_cc^3}{4G\hbar}
=\frac{\pi R_c^2c^3}{G\hbar}
\]

and the canonically normalized horizontal causal-diamond temperature

\[
k_BT_c=\frac{\hbar c}{2\pi R_c}.
\]

Then

\[
\boxed{
k_BT_c\frac{S_c}{k_B}=\frac{c^4R_c}{2G}=E_{\rm MS,c}.}
\]

The flat Friedmann relation gives

\[
E_{\rm MS,c}=\rho_{\rm crit,c}V_c.
\]

Thus

\[
\boxed{
\frac{k_BT_c}{V_c}\frac{S_c}{k_B}=\rho_{\rm crit,c}.
}
\]

This is the dimensional bridge. The BKM metric supplies dimensionless
response; horizon temperature supplies energy per information unit; the
causal volume supplies density; Friedmann marginality identifies the
result with the critical density.

\begin{figure}
\centering
\includegraphics[width=0.96\textwidth,height=\textheight]{figures/scale_capacity_closure.pdf}
\caption{The homogeneous amplitude closes by composing binary BKM
geometry, modular free energy, and horizon thermodynamics.}
\end{figure}

\subsection{Closed amplitude}\label{closed-amplitude}

Using

\[
G^{\perp}_{NN}(N)
=\mathfrak R_c\frac{S_c}{k_B}\operatorname{sech}^2(N-N_c),
\]

the source becomes

\[
\boxed{
\rho_X(N)=\frac{\mathfrak R_c}{2}\rho_{\rm crit,c}\operatorname{sech}^2(N-N_c).
}
\]

For the fundamental value \(\mathfrak R_c=1\),

\[
\boxed{
\rho_X(N)=\frac12\rho_{\rm crit,c}\operatorname{sech}^2(N-N_c).
}
\]

At the crossing,

\[
\Omega_{X,c}=\frac12.
\]

Spatial flatness implies

\[
\rho_{\rm ordinary,c}=\rho_{\rm crit,c}-\rho_{X,c}=\frac12\rho_{\rm crit,c},
\]

and hence

\[
\boxed{
\rho_X(N_c)=\rho_{\rm ordinary}(N_c).
}
\]

The former postulate \(r_c=1\) is thereby derived from the
scale-capacity law. The exact equality includes radiation and any other
non-dark contribution.

\subsection{General spatial dimension}\label{general-spatial-dimension}

Let \(d\) be the number of spatial dimensions. The apparent-horizon area
and volume scale as

\[
A=\Omega_{d-1}R^{d-1},
\qquad
V=\frac{\Omega_{d-1}}{d}R^d.
\]

Einstein-Friedmann geometry gives

\[
\frac{k_BT_c(S_c/k_B)}{V_c}
=\frac{2}{d-1}\rho_{\rm crit,c}.
\]

Therefore

\[
\boxed{
\Omega_{X,c}=\frac{\mathfrak R_c}{d-1},
}
\]

and

\[
\boxed{
\frac{\rho_{X,c}}{\rho_{\rm ordinary,c}}
=\frac{\mathfrak R_c}{d-1-\mathfrak R_c}.
}
\]

For \(\mathfrak R_c=1\),

\[
\frac{\rho_{X,c}}{\rho_{\rm ordinary,c}}=\frac1{d-2}.
\]

Equal ordinary and response densities occur in three spatial dimensions:

\[
\boxed{
\mathfrak R_c=1,
\quad
\rho_X=\rho_{\rm ordinary}
\quad\Longleftrightarrow\quad
d=3.
}
\]

This does not independently prove that space must have three dimensions.
It shows that the unit scale-capacity law and equal self-dual partition
are mutually compatible precisely in the observed dimension.

\begin{figure}
\centering
\includegraphics[width=0.78\textwidth,height=\textheight]{figures/dimension_crossing_ratio.pdf}
\caption{The dimension dependence of the self-dual crossing.}
\end{figure}

\subsection{Status of the closure}\label{status-of-the-closure}

The closure contains one physical law:

\[
\boxed{
\left.
\frac{k_B}{S_c}G^{\perp}_{NN}
\right|_{N_c}=1.
}
\]

The supporting evidence is:

\begin{itemize}
\tightlist
\item
  capacity equals entropy in controlled spherical Einstein-holographic
  settings ({[}26{]});
\item
  modular-Hamiltonian variance equal to horizon entropy has been
  proposed for flat, de Sitter, and suitable holographic causal diamonds
  ({[}27,28{]});
\item
  the BKM Hessian equals gravitational canonical energy in controlled
  holographic perturbation theory ({[}12{]});
\item
  apparent-horizon thermodynamics gives the exact Hawking--Friedmann
  dimensional conversion ({[}24,25{]}).
\end{itemize}

The law has not been proved for a dynamical FLRW causal wall. It is the
theory's equivalence principle and its most direct target for
independent derivation or falsification.

\section{Exact homogeneous dynamics and the sign of
acceleration}\label{exact-homogeneous-dynamics-and-the-sign-of-acceleration}

\subsection{Flat FLRW scale
dictionary}\label{flat-flrw-scale-dictionary}

For a spatially flat FLRW metric

\[
\mathrm ds^2=a^2(\eta)(-\mathrm d\eta^2+\mathrm d\mathbf x^2),
\]

choose the flat conformal representative. Then

\[
\sigma=\frac1a,
\qquad
N=\ln a=-\ln\sigma.
\]

The following identities are exact:

\[
\boxed{\sigma'=-H,}
\]

\[
\boxed{\mathcal R=-\frac{\sigma}{\sigma'},}
\]

\[
\boxed{\mathcal R'=q,}
\]

and

\[
\boxed{
q=-1+\frac{\sigma\sigma''}{\sigma'^2}
=\frac{(\ln\sigma)''}{[(\ln\sigma)']^2}.
}
\]

Here a prime denotes conformal-time differentiation and
\(\mathcal R=1/(aH)\) is the comoving Hubble radius.

The homogeneous null-energy condition is

\[
\rho+p\ge0
\quad\Longleftrightarrow\quad
\sigma''\ge0,
\]

so it is convexity of the scale section. Accelerated expansion is

\[
q<0
\quad\Longleftrightarrow\quad
(\ln\sigma)''<0,
\]

so it is log-concavity of scale.

The tractor norm becomes

\[
\boxed{
I^2=\frac12\sigma\sigma''-\sigma'^2
=-\frac12(1-q)H^2.
}
\]

\subsection{Closed background
equation}\label{closed-background-equation}

With ordinary radiation and matter plus the scale-capacity source,

\[
\boxed{
H^2(N)=\frac{8\pi G}{3c^2}
\left[
\rho_{m0}e^{-3N}
+\rho_{r0}e^{-4N}
+\frac12\rho_{\rm crit,c}\operatorname{sech}^2(N-N_c)
\right].
}
\]

At the self-dual point,

\[
\frac12\rho_{\rm crit,c}
=\rho_m(N_c)+\rho_r(N_c)+\cdots.
\]

Present flatness determines \(N_c\) from

\[
\boxed{
1=\Omega_{m0}+\Omega_{r0}
+\left[
\Omega_{m0}e^{-3N_c}
+\Omega_{r0}e^{-4N_c}
\right]\operatorname{sech}^2N_c.
}
\]

Thus \(N_c\) is not an independently fitted dark parameter. Its
intrinsic meaning is fixed by self-duality; its position relative to
today is fixed by the measured ordinary state and flatness.

\subsection{Equation of state and structural
invariant}\label{equation-of-state-and-structural-invariant}

The density law gives

\[
\frac{\mathrm d\ln\rho_X}{\mathrm dN}
=-2\tanh(N-N_c).
\]

Separate conservation gives

\[
\boxed{
w_X(N)=-1+\frac23\tanh(N-N_c).
}
\]

Let

\[
X:=1+w_X.
\]

Then

\[
\boxed{
X'=\frac23-\frac32X^2.
}
\]

The two fixed points are

\[
w_-=-\frac53,
\qquad
w_+=-\frac13.
\]

The observed history is the unique heteroclinic orbit from the early to
the late fixed point. The one density maximum and one \(w_X=-1\)
crossing are not separately adjusted features.

The binary normalization becomes the differential invariant

\[
\boxed{
9(1+w_X)^2+6w_X'=4.
}
\]

This relation is independent of the amplitude and crossing date. It is
the primary structural test of the theory.

\subsection{Why acceleration begins and
ends}\label{why-acceleration-begins-and-ends}

The total deceleration parameter is

\[
q=-1+\frac32
\frac{
\rho_m+\frac43\rho_r+(1+w_X)\rho_X
}{
\rho_m+\rho_r+\rho_X
}.
\]

Acceleration occurs when

\[
\boxed{
(2-3[1+w_X])\rho_X>\rho_m+2\rho_r.
}
\]

The response is negligible in the remote past, rises through the
ordinary budget, and later decays as \(a^{-2}\). Consequently the
inequality is satisfied only on a finite interval.

The model therefore predicts two zeroes of \(q\):

\begin{enumerate}
\def\labelenumi{\arabic{enumi}.}
\tightlist
\item
  an observed past transition from deceleration to acceleration;
\item
  a future transition back to nonacceleration.
\end{enumerate}

The phrase ``the sign flip of cosmic acceleration'' ordinarily refers to
the first. The second is an independent future prediction.

\subsection{Benchmark}\label{benchmark}

For

\[
\Omega_{m0}=0.310598,
\qquad
\Omega_{r0}=9.15\times10^{-5},
\]

the closed solution is

\[
\boxed{
N_c=-0.2940066,
\qquad
z_c=0.3417927.
}
\]

The crossing density in present critical units is

\[
\frac{\rho_*}{\rho_{\rm crit,0}}=0.7506311.
\]

The exact ordinary-sector equality is

\[
\frac{\rho_*}{\rho_m(N_c)+\rho_r(N_c)}=1.
\]

Relative to dust alone,

\[
\frac{\rho_*}{\rho_m(N_c)}=1.0003953.
\]

The present and transition observables are

\begin{longtable}[]{@{}lr@{}}
\toprule\noalign{}
quantity & prediction \\
\midrule\noalign{}
\endhead
\bottomrule\noalign{}
\endlastfoot
\(w_0\) & \(-0.809454\) \\
tangent \(w_a\) & \(-0.612205\) \\
\(q_0\) & \(-0.336902\) \\
\(j_0\) & \(-0.111246\) \\
\(q(N_c)\) & \(-0.249901\) \\
horizon index \(\mu_A(N_c)\) & \(0.624951\) \\
acceleration entry & \(z=0.785694\) \\
acceleration exit & \(a/a_0=11.7865\) \\
\end{longtable}

\begin{figure}
\centering
\includegraphics[width=0.94\textwidth,height=\textheight]{figures/rigid_history_v7.pdf}
\caption{The capacity-normalized background. The self-dual crossing, the
observed acceleration entry, and the predicted exit are distinct
events.}
\end{figure}

\subsection{Future causal character}\label{future-causal-character}

At late times,

\[
\rho_X\propto a^{-2},
\qquad
w_X\to-\frac13.
\]

If the residual scalar floor is exactly zero,

\[
H\propto a^{-1},
\qquad
\dot a\to\text{constant},
\qquad
a(t)\sim t.
\]

The future conformal-time integral diverges,

\[
\int^\infty\frac{\mathrm dt}{a(t)}\sim\int^\infty\frac{\mathrm dt}{t}=\infty,
\]

so there is no permanent future event horizon.

A positive residual \(\Lambda_g\), however small today, eventually
dominates and changes the asymptotic state to de Sitter. Exact zero and
observationally negligible are therefore distinct global sectors.

\section{Local gravitational sourcing and the vacuum
catastrophe}\label{local-gravitational-sourcing-and-the-vacuum-catastrophe}

\subsection{One causal boost charge {[}CONJECTURAL UNITY
PRINCIPLE{]}}\label{one-causal-boost-charge-conjectural-unity-principle}

For an equilibrium causal cut, propose that the modular boost charge of
the restricted state and the gravitational Noether boost charge of the
cut are two representations of one causal charge.

The state side gives the entanglement first law ({[}7,9,10{]})

\[
\delta S_{\rm out}=\delta\left\langle K\right\rangle.
\]

The geometric side gives

\[
S_{\rm grav}=\frac{k_BA}{4\ell_P^2}
\]

in the Einstein regime. Raychaudhuri focusing relates area variation to
\(R_{ab}k^ak^b\). Equating the two boost-charge responses for every
local null direction yields

\[
R_{ab}k^ak^b=\frac{8\pi G}{c^4}T_{ab}k^ak^b.
\]

A Lorentzian tensor lemma then gives

\[
R_{ab}^{\circ}=\frac{8\pi G}{c^4}T_{ab}^{\circ},
\]

and therefore the tractor transport equation in Section 4.

This is a conditional derivation combining established ingredients with
one unity principle. It is not derived from the binary late-time model.

\subsection{Vacuum offset as a vertical/central
direction}\label{vacuum-offset-as-a-verticalcentral-direction}

A constant Hamiltonian shift leaves the normalized state unchanged:

\[
\frac{e^{-\beta(H+C\mathbf1)}}{\operatorname{tr}e^{-\beta(H+C\mathbf1)}}
=\frac{e^{-\beta H}}{\operatorname{tr}e^{-\beta H}}.
\]

It therefore has zero BKM length. The corresponding stress shift is
annihilated by the trace-free source equation.

The proposed retyping is:

\[
\boxed{
\begin{aligned}
\text{central vacuum offset}
&=\text{vertical normalization direction},\\
\text{physical state response}
&=\text{horizontal BKM displacement},\\
\text{cosmological scalar curvature}
&=\text{global lift data}.
\end{aligned}}
\]

This shows why the local causal dynamics need not respond to the QFT
zero of energy. It does not determine or stabilize the global lift.

\subsection{One explicit global
completion}\label{one-explicit-global-completion}

Manifestly local vacuum-energy sequestering introduces four-form sectors
({[}26,27{]}) whose equations force the cosmological variables to be
spacetime constants while global flux constraints subtract the
spacetime-averaged vacuum contribution. Schematically,

\[
M^2G^a{}_b
=T^a{}_b-\frac14\delta^a_b\left\langle T\right\rangle-\Delta\Lambda\,\delta^a_b.
\]

Writing

\[
T^a{}_b=-V_{\rm vac}\delta^a_b+\tau^a{}_b
\]

gives

\[
M^2G^a{}_b
=\tau^a{}_b-\frac14\delta^a_b\left\langle \tau\right\rangle-\Delta\Lambda\,\delta^a_b.
\]

The spacetime-filling matter-loop vacuum cancels from the source. A
Gauss--Bonnet extension has been proposed to protect the residual
against the cutoff-dominated graviton-loop sector within effective field
theory.

This is one published completion, not a consequence of modular--Weyl
geometry. The master programme's claim is that the local quotient and
global top-form lift have compatible mathematical types.

\subsection{What remains of the cosmological
constant}\label{what-remains-of-the-cosmological-constant}

The surviving residual

\[
\Lambda_{\rm res}
\]

is a global flux or superselection datum. The rigid branch takes

\[
\Lambda_{\rm res}=0
\]

exactly. A small positive value is not equivalent: it eventually
dominates and restores a future event horizon.

The programme does not yet derive which flux sector nature selects.

\subsection{Relation to the scale-capacity
source}\label{relation-to-the-scale-capacity-source}

The scale-capacity source is automatically blind to a central modular
shift because its primitive is a covariance:

\[
G^{\mathrm{BKM}}_{NN}\sim\operatorname{Var}(K_N),
\]

and

\[
\operatorname{Var}(K+\alpha\mathbf1)=\operatorname{Var}(K).
\]

This does not mean that every quantum contribution proportional to the
metric has been removed from a fully renormalized gravitational
effective action. It means that the \textbf{local horizontal response
law} depends only on noncentral state distinguishability. The
cosmological constant and radiative stability of the scalar lift are
logically separate.

The theory therefore assigns the two problems to different mathematical
slots:

\[
\boxed{
\begin{aligned}
\text{transient late response}
&\longleftrightarrow
\text{horizontal BKM free-energy curvature},\\
\text{constant vacuum offset}
&\longleftrightarrow
\text{central/global norm sector}.
\end{aligned}}
\]

This separation is one of the principal motivations for the framework.
It does not replace the need to specify the global sector.

\section{Horizons, black holes, entropy, and the meanings of
time}\label{horizons-black-holes-entropy-and-the-meanings-of-time}

\subsection{Four distinct notions of
time}\label{four-distinct-notions-of-time}

The framework uses several ordered parameters that must not be
conflated.

\begin{longtable}[]{@{}
  >{\raggedright\arraybackslash}p{(\columnwidth - 2\tabcolsep) * \real{0.5000}}
  >{\raggedright\arraybackslash}p{(\columnwidth - 2\tabcolsep) * \real{0.5000}}@{}}
\toprule\noalign{}
\begin{minipage}[b]{\linewidth}\raggedright
Symbol or notion
\end{minipage} & \begin{minipage}[b]{\linewidth}\raggedright
Meaning
\end{minipage} \\
\midrule\noalign{}
\endhead
\bottomrule\noalign{}
\endlastfoot
causal order & a partial order on events; no metric duration is
implied \\
proper or clock time \(\tau\) & metric length along a timelike
worldline \\
Weyl scale time \(N\) & logarithmic change of the scale section,
\(N=\ln(a/a_c)\) \\
vertical modular parameter \(s\) & automorphism parameter within one
algebra/state fiber \\
horizontal state coordinate \(\theta\) & relative modular polarization
across a family of states \\
\end{longtable}

The v5.0 inconsistency arose from treating the last two as the same
parameter. They are different directions even when they use the same
normal boost generator.

\subsection{Running horizon index}\label{running-horizon-index}

For a spherical horizon of areal radius \(R_H\), define

\[
\boxed{
\mu_H:=\frac{|\kappa_H|R_H}{c^2}.
}
\]

It compares a normal-boost rate with an inverse geometric scale. It is a
state-dependent index, not the Ruble number and not a measure of a total
``amount of gravity.''

For a flat FLRW apparent horizon,

\[
R_A=\frac{c}{H},
\]

and the Kodama--Hayward surface gravity gives

\[
\boxed{
\mu_A=\frac{1-q}{2}=-\frac{I^2}{H^2}.
}
\]

The first equality is horizon kinematics; the second is scale-tractor
geometry. In four-dimensional spherical Einstein gravity,

\[
S_A=k_B\frac{\pi R_A^2c^3}{G\hbar},
\qquad
E_A=\frac{c^4R_A}{2G},
\]

so, with the positive Kodama--Hayward temperature,

\[
\mu_A=\frac{T_AS_A}{E_A}.
\]

Representative values are

\begin{longtable}[]{@{}lrrr@{}}
\toprule\noalign{}
regime & \(w\) & \(q\) & \(\mu_A\) \\
\midrule\noalign{}
\endhead
\bottomrule\noalign{}
\endlastfoot
radiation & \(1/3\) & \(1\) & \(0\) \\
matter & \(0\) & \(1/2\) & \(1/4\) \\
coasting / acceleration threshold & \(-1/3\) & \(0\) & \(1/2\) \\
de Sitter & \(-1\) & \(-1\) & \(1\) \\
\end{longtable}

Therefore

\[
\boxed{
\ddot a>0\iff\mu_A>\frac12.
}
\]

The value \(1/2\) also equals \(|\kappa|R/c^2\) for a four-dimensional
Schwarzschild horizon. This identifies a shared dimensionless
surface-gravity balance; it does not identify an FLRW spacetime with a
Schwarzschild spacetime.

\subsection{Clock allocation}\label{clock-allocation}

The dimensionless apparent-horizon entropy is

\[
\mathcal S_A=\frac{S_A}{k_B}=\frac{\pi c^5}{G\hbar H^2}.
\]

Hence

\[
\frac{\mathrm d\ln\mathcal S_A}{\mathrm dN}=4(1-\mu_A).
\]

Define a geometrically normalized vertical horizon rapidity potential
\(\eta_A\) by

\[
\frac{\mathrm d\eta_A}{\mathrm dt}=\frac{|\kappa_A|}{c}.
\]

Since \(\mathrm dN/\mathrm dt=H\),

\[
\frac{\mathrm d\eta_A}{\mathrm dN}=\mu_A.
\]

Thus

\[
\boxed{
\mathrm dN=\mathrm d\eta_A+\frac14\mathrm d\ln\mathcal S_A.
}
\]

One Weyl e-fold decomposes exactly into vertical horizon-rapidity
advance and horizon-information growth. The horizontal fundamental law
is instead

\[
\frac{\mathrm d\theta}{\mathrm dN}=1.
\]

Consequently,

\[
\boxed{
\theta'-\eta_A'
=1-\mu_A
=\frac14(\ln\mathcal S_A)'.
}
\]

This is the corrected relation among horizontal state motion, vertical
horizon motion, and information capacity.

\subsection{Smarr and Hawking relations as scale
homogeneity}\label{smarr-and-hawking-relations-as-scale-homogeneity}

For a one-scale asymptotically flat black hole in \(D\) spacetime
dimensions,

\[
E\propto R_H^{D-3},
\qquad
S\propto R_H^{D-2}.
\]

Euler homogeneity and \(\mathrm dE=T\mathrm dS\) give

\[
\boxed{
(D-3)E=(D-2)TS.
}
\]

The Smarr relation is therefore the compatibility of thermodynamic
variation with scale homogeneity. In four dimensions, \(E=2TS\).

The cosmological relation used in the amplitude closure is not the
stationary Schwarzschild Smarr relation. It uses the canonically
normalized horizontal causal-diamond temperature,

\[
k_BT_c=\frac{\hbar c}{2\pi R_c},
\]

together with the apparent-horizon area law and Misner--Sharp
marginality to obtain

\[
T_cS_c=E_{\rm MS,c}.
\]

This distinction between vertical dynamical surface gravity and
horizontal modular normalization is essential.

\subsection{Holography as a controlled
laboratory}\label{holography-as-a-controlled-laboratory}

In a holographic code subspace, the JLMS relation takes the schematic
form

\[
K_A^{\rm CFT}
=\frac{\widehat A(\chi_A)}{4G_N\hbar}
+K_a^{\rm bulk}
+\text{central term}
+O(G_N).
\]

Boundary modular charge decomposes into a geometric area charge and a
bulk modular charge. In related controlled settings:

\begin{itemize}
\tightlist
\item
  the Hessian of relative entropy equals gravitational canonical energy;
\item
  modular Berry curvature maps to gravitational symplectic structure;
\item
  crossed-product gravitational algebras make generalized entropy a
  genuine von Neumann entropy for suitable subregions.
\end{itemize}

These results do not prove the FLRW scale-capacity law. They establish
that modular charge, area, information Hessian, symplectic form, and
gravity can form one coherent subregion structure.

\section{Horizontal perturbation geometry and the open spacetime
lift}\label{horizontal-perturbation-geometry-and-the-open-spacetime-lift}

The homogeneous source is closed by the scale-capacity law. A complete
physical theory must also determine inhomogeneous perturbations. This
section records what is fixed by the binary state geometry and what
remains open.

\subsection{Canonical Witten--Darboux
pair}\label{canonical-wittendarboux-pair}

The binary information potential is

\[
\Psi(\theta)=\ln(2\cosh\theta),
\]

with

\[
\eta(\theta)=\Psi'(\theta)=\tanh\theta,
\qquad
\eta'(\theta)=\operatorname{sech}^2\theta.
\]

Define first-order operators

\[
\mathcal A=\partial_\theta+\eta,
\qquad
\mathcal A^\dagger=-\partial_\theta+\eta.
\]

Their partner Hamiltonians are

\[
\mathcal H_-:=\mathcal A^\dagger\mathcal A
=-\partial_\theta^2+\eta^2-\eta',
\]

\[
\mathcal H_+:=\mathcal A\mathcal A^\dagger
=-\partial_\theta^2+\eta^2+\eta'.
\]

Using \(\eta^2+\eta'=1\) gives the exact factorization

\[
\boxed{
\mathcal H_-=-\partial_\theta^2+1-2\operatorname{sech}^2\theta,
}
\]

\[
\boxed{
\mathcal H_+=-\partial_\theta^2+1.
}
\]

Equivalently, define the two-component Dirac/Witten operator

\[
\boxed{
\mathcal D_\Psi=
\begin{pmatrix}
0&\mathcal A^\dagger\\
\mathcal A&0
\end{pmatrix},
\qquad
\mathcal D_\Psi^2=
\begin{pmatrix}
\mathcal H_-&0\\
0&\mathcal H_+
\end{pmatrix}.
}
\]

This pair is not an independently guessed Pöschl--Teller ansatz. It is
generated canonically by the binary log-partition potential;
reflectionless \(\operatorname{sech}^2\) operators are a standard
exactly solvable class ({[}33{]}).

\subsection{Bound mode, BKM density, and
transparency}\label{bound-mode-bkm-density-and-transparency}

The zero-mode equation

\[
\mathcal A\psi_0=0
\]

has normalized solution

\[
\boxed{
\psi_0(\theta)=\frac1{\sqrt2}\operatorname{sech}\theta.
}
\]

Therefore

\[
\boxed{
2|\psi_0|^2=\operatorname{sech}^2\theta=G^{\mathrm{BKM}}_{\theta\theta}.
}
\]

The same function is simultaneously:

\begin{itemize}
\tightlist
\item
  the BKM/Fisher metric density;
\item
  the normalized dark-history shape;
\item
  the density of the unique bound mode of the Witten pair.
\end{itemize}

For continuum energy \(1+k^2\), applying \(\mathcal A^\dagger\) to a
free wave gives

\[
\psi_k(\theta)=(-ik+\tanh\theta)e^{ik\theta}.
\]

At both asymptotic ends there is no \(e^{-ik\theta}\) component, so

\[
\boxed{R(k)=0.}
\]

The continuum is transmitted with a phase shift; it is not absent. The
correct statement is one normalizable bound mode plus a reflectionless
continuum.

The Witten index is one, matching the unit topological charge of the
Fisher kink. The total Fisher length is \(\pi\), and the one-bound-state
Levinson phase is also \(\pi\); these are consequences of the same
factorization rather than independent evidence.

\subsection{What this pair does and does not
prove}\label{what-this-pair-does-and-does-not-prove}

The Witten pair closes the \textbf{internal horizontal operator
geometry}. It explains why the fundamental pair is \(\ell=1\), why there
is one localized mode, and why the continuum is reflectionless.

It does not yet identify these operators with the physical scalar,
vector, or tensor perturbation operators of an FLRW spacetime. A
spacetime perturbation theory must specify:

\begin{itemize}
\tightlist
\item
  how inhomogeneous deformations of causal cuts induce horizontal state
  tangents;
\item
  how the two-component Witten complex is embedded in spacetime
  constraints;
\item
  how spatial gradients and causal propagation enter;
\item
  which observable combination couples to matter and metric
  perturbations;
\item
  whether the resulting system is ghost-free, hyperbolic, and regular at
  the crossing.
\end{itemize}

Thus the pair completion is canonical, but the spacetime lift remains
open.

\subsection{Negative result: ordinary matter growth is not the pair
operator}\label{negative-result-ordinary-matter-growth-is-not-the-pair-operator}

The standard smooth-dark-response matter-growth equation can be reduced
to a Schrödinger-like zero-energy equation, but its effective potential
contains ordinary matter and Hubble-friction terms. It is not
proportional to \(\rho_X\) and is not the Pöschl--Teller operator above.

Therefore transparency must not be imposed on the ordinary matter-growth
equation. The Witten pair belongs to the internal horizontal response
sector, not to a rebranding of standard growth.

\subsection{Negative result: canonical single-field
completion}\label{negative-result-canonical-single-field-completion}

For a canonical single scalar, the kinetic normalization is proportional
to \(1+w_X\). In the rigid history,

\[
1+w_X=\frac23\tanh\theta,
\]

which is negative throughout the entire pre-crossing branch. At the
crossing, the corresponding single-field mode function behaves as

\[
z\sim|\theta|^{1/2},
\]

so

\[
\frac{z''}{z}\sim-\frac1{4\theta^2},
\]

at the critical inverse-square coupling. The horizontal state coordinate
therefore cannot be interpreted as an ordinary canonical local scalar
field.

This is not a defect of the homogeneous free-energy construction. It
says that the physical completion must be collective, constrained,
multi-component, algebraic, or otherwise noncanonical.

\subsection{Perturbative stopping
condition}\label{perturbative-stopping-condition}

A perturbative completion is adequate when it derives, without arbitrary
functions,

\begin{enumerate}
\def\labelenumi{\arabic{enumi}.}
\tightlist
\item
  a conserved \(\delta T^X_{ab}\);
\item
  a regular crossing;
\item
  a finite number of propagating or constrained modes fixed by the
  Witten pair;
\item
  a definite gradient structure and sound/cone speed;
\item
  CMB, lensing, and growth responses that can be calculated in a
  Boltzmann code.
\end{enumerate}

Until then, the background equations are closed but the theory is not a
complete cosmological perturbation theory.

\section{Quantitative predictions, observational status, and kill
conditions}\label{quantitative-predictions-observational-status-and-kill-conditions}

\subsection{Parameter-free rigid
branch}\label{parameter-free-rigid-branch}

The rigid branch uses

\[
\mathfrak R_c=1,
\qquad
\varrho_\perp=1,
\qquad
\Lambda_{\rm res}=0,
\qquad
k=0.
\]

The present ordinary matter and radiation abundances are observational
inputs, not dark-sector fit parameters. The crossing date is then fixed
by present flatness.

\subsection{P1 --- differential shape
invariant}\label{p1-differential-shape-invariant}

The primary test is

\[
\boxed{
9(1+w_X(z))^2+6\frac{\mathrm dw_X}{\mathrm dN}=4.
}
\]

This constrains the whole reconstructed history. It is independent of
the amplitude and crossing date.

\textbf{Kill condition K1.} A statistically significant redshift
dependence of the left-hand side, after accounting for reconstruction
covariance, rejects the binary affine theory.

\subsection{P2 --- local tangent locus}\label{p2-local-tangent-locus}

At the present epoch, the exact history has CPL tangent

\[
\boxed{
w_a=\frac32(1+w_0)^2-\frac23.
}
\]

For the benchmark,

\[
(w_0,w_a)=(-0.80945,-0.61221).
\]

This is a compressed diagnostic only; the exact pulse should be fitted
directly.

\textbf{Kill condition K2.} A clean posterior excluding this locus
rejects the rigid branch.

\subsection{P3 --- self-dual crossing and
equality}\label{p3-self-dual-crossing-and-equality}

The model predicts

\[
\boxed{
z_c=0.34179,}
\]

and the exact coincidence

\[
\boxed{
\rho_X(N_c)=\rho_{\rm ordinary}(N_c),
\qquad
w_X(N_c)=-1.
}
\]

Relative to dust alone, the ratio is \(1.000395\) because radiation is
small but nonzero.

\textbf{Kill condition K3.} Independent reconstructions of the \(w=-1\)
crossing and ordinary--dark equality that disagree beyond errors reject
the scale-capacity normalization.

\subsection{P4 --- observed acceleration entry and future
exit}\label{p4-observed-acceleration-entry-and-future-exit}

The benchmark predicts

\[
\boxed{z_{\rm acc}=0.78569}
\]

for the past entry into acceleration, and

\[
\boxed{a_{\rm exit}/a_0=11.7865}
\]

for the future exit. The crossing/equality event is not the acceleration
onset.

The future tends to

\[
w_X\to-\frac13,
\qquad
a(t)\sim t,
\]

with no permanent event horizon if and only if the residual floor is
exactly zero.

\subsection{P5 --- present cosmography}\label{p5-present-cosmography}

The rigid benchmark predicts

\[
\boxed{q_0=-0.33690,}
\]

\[
\boxed{j_0=-0.11125.}
\]

The jerk is separated by order unity from the flat-\(\Lambda\)CDM value
\(j=1\).

\textbf{Kill condition K4.} A robust model-independent reconstruction
near \(j_0=1\) and inconsistent with the value above strongly rejects
the rigid history.

\subsection{P6 --- dimension and crossing
fraction}\label{p6-dimension-and-crossing-fraction}

In \(d\) spatial dimensions, the scale-capacity law predicts

\[
\boxed{
\Omega_{X,c}=\frac{1}{d-1},
\qquad
\frac{\rho_{X,c}}{\rho_{\rm ordinary,c}}=\frac1{d-2}.
}
\]

The equal partition occurs only for \(d=3\). This is not an empirical
prediction within our already observed dimension; it is a structural
consistency condition relevant to dimensional-selection attempts.

\subsection{P7 --- horizontal operator
structure}\label{p7-horizontal-operator-structure}

The internal response geometry predicts the Witten pair

\[
\mathcal H_-=-\partial_\theta^2+1-2\operatorname{sech}^2\theta,
\qquad
\mathcal H_+=-\partial_\theta^2+1,
\]

with one bound mode and reflectionless continuum.

\textbf{Kill condition K5.} A derived spacetime completion whose
horizontal response has a different bound-state count or cannot realize
the pair invalidates the claimed canonical perturbation structure.

\subsection{Current background
comparison}\label{current-background-comparison}

The previously completed background-only analysis of public DESI
distance data and Pantheon+ supernovae gives approximately
({[}34,36,37{]})

\begin{longtable}[]{@{}lrr@{}}
\toprule\noalign{}
model & background shape parameters & joint \(\chi^2\) \\
\midrule\noalign{}
\endhead
\bottomrule\noalign{}
\endlastfoot
flat \(\Lambda\)CDM & 1 & 1401.63 \\
rigid capacity-normalized history & 1 & 1398.29 \\
flexible pulse & 3 & 1397.26 \\
\end{longtable}

The rigid history is viable and modestly improves this restricted
statistic at equal background shape-parameter count. These values are
not a discovery significance. The comparison excludes primary CMB, full
lensing, growth, and a derived dark perturbation likelihood.

\subsection{Neutrino residual}\label{neutrino-residual}

Current cosmological analyses infer an unusually tight neutrino-mass
bound in flat \(\Lambda\)CDM and, under a formal continuation, can
prefer an effective negative response ({[}35{]}). Allowing evolving
late-time geometry substantially weakens that result.

The rigid history has the correct qualitative structure to address the
residual:

\begin{itemize}
\tightlist
\item
  it modifies low- and intermediate-redshift expansion with opposite
  signs relative to a matched \(\Lambda\)CDM history;
\item
  its smooth-growth response increases CMB lensing in the direction
  opposite positive neutrino mass;
\item
  a geometry-level pushforward places weight near a normal-ordering sum
  of order \(0.07\)--\(0.08\) eV.
\end{itemize}

However, null ensembles show that the lensing anti-alignment is common
among smooth positive transient histories. The neutrino calculation is
therefore a necessary class-membership test, not a distinctive
confirmation.

A direct CLASS/CAMB implementation with a derived perturbation closure
is required before the neutrino anomaly can be claimed as a prediction
or resolution.

\subsection{Epistemic status of the observational
evidence}\label{epistemic-status-of-the-observational-evidence}

The current evidence supports:

\begin{enumerate}
\def\labelenumi{\arabic{enumi}.}
\tightlist
\item
  compatibility of the rigid background with present distance data;
\item
  current dynamical-dark-energy summaries lying near the invariant's
  natural value;
\item
  a redshift structure capable of reducing the neutrino residual.
\end{enumerate}

It does not establish:

\begin{enumerate}
\def\labelenumi{\arabic{enumi}.}
\tightlist
\item
  the scale-capacity principle for FLRW causal walls;
\item
  a complete perturbation theory;
\item
  the global zero-floor sector;
\item
  a statistically preferred particle-mass posterior under the full
  model.
\end{enumerate}

\section{Economy, epicycles, and the novelty
test}\label{economy-epicycles-and-the-novelty-test}

\subsection{Four distinct economies}\label{four-distinct-economies}

Parameter count alone is not explanatory economy. Distinguish:

\[
\begin{aligned}
\mathcal E_{\rm fit}&=\text{number of continuously fitted parameters},\\
\mathcal E_{\rm law}&=\text{number of independently postulated laws or functions},\\
\mathcal E_{\rm ontology}&=\text{number of new fields, fluids, or forces},\\
\mathcal E_{\rm consequence}&=\text{number of predictions locked together}.
\end{aligned}
\]

Epicycles fail principally because corrections are independently
adjustable. A fixed but unexplained function can have zero fitted
parameters and still be ad hoc.

The rigid causal-scale theory has no continuously fitted dark-history
shape parameter once its equivalence principle and global sector are
fixed. More importantly, one structure ties together:

\begin{itemize}
\tightlist
\item
  one density maximum;
\item
  one \(w=-1\) crossing;
\item
  the full differential invariant;
\item
  one acceleration interval;
\item
  a future separatrix;
\item
  the self-dual amplitude;
\item
  a canonical horizontal Witten pair;
\item
  local blindness to central vacuum shifts.
\end{itemize}

This is a stronger kind of economy than low parameter count alone.

\subsection{Economical competitors}\label{economical-competitors}

Several existing cosmologies are also economical at the background level
({[}38-40{]}):

\begin{itemize}
\tightlist
\item
  flat \(\Lambda\)CDM has one constant vacuum term;
\item
  phenomenological emergent dark-energy models can have no extra fitted
  parameter;
\item
  vacuum metamorphosis derives a late transition from a
  curvature-triggered quantum effect;
\item
  running-vacuum models use a small number of QFT-inspired coefficients;
\item
  hyperbolic-tangent crossing models can closely resemble the rigid
  background with two fitted shape parameters.
\end{itemize}

The present framework is not unique because it uses a
\(\operatorname{sech}^2\) profile or because it fits without a free
\(w_0,w_a\) pair. Its claim to novelty is the proposed identity

\[
\boxed{
\text{Weyl-scale translation}
\longrightarrow
\text{horizontal modular-state displacement}
\longrightarrow
\text{BKM capacity}
\longrightarrow
\text{gravitational free-energy density}.
}
\]

If that chain is independently constructed, the profile is a
consequence. If the chain is only inferred from the fitted expansion
history, the model is an effective fluid in new notation.

\subsection{Musical-chairs test}\label{musical-chairs-test}

The decisive criterion is:

\[
\boxed{
\begin{array}{ll}
\theta\text{ is defined only by rearranging }H(z)
&\Rightarrow\text{relabelled fluid},\\[1mm]
\theta\text{ is defined by an independent causal-state construction}
&\Rightarrow\text{new physical structure}.
\end{array}}
\]

The present theory addresses this by defining \(\theta\) as the
noncentral Connes-cocycle coordinate of a scale-indexed wall-state
family and by identifying the active quotient with normal chirality. The
full construction for a dynamical FLRW wall remains to be carried out.

\subsection{Assumption ledger after
closure}\label{assumption-ledger-after-closure}

The homogeneous theory now contains:

\begin{longtable}[]{@{}
  >{\raggedright\arraybackslash}p{(\columnwidth - 2\tabcolsep) * \real{0.5000}}
  >{\raggedright\arraybackslash}p{(\columnwidth - 2\tabcolsep) * \real{0.5000}}@{}}
\toprule\noalign{}
\begin{minipage}[b]{\linewidth}\raggedright
item
\end{minipage} & \begin{minipage}[b]{\linewidth}\raggedright
status
\end{minipage} \\
\midrule\noalign{}
\endhead
\bottomrule\noalign{}
\endlastfoot
causal conformal geometry plus scale & standard mathematical
factorization \\
local modular/gravitational boost-charge identity & physical equivalence
principle used for the local source equation \\
fundamental normal chirality quotient & structural identification of the
active homogeneous mode \\
affine cocycle soldering & theorem given rank-one ratio dependence and
measurability \\
fundamental representation \(\varrho_\perp=1\) & physical representation
choice / unit character \\
scale--capacity number \(\mathfrak R_c=1\) & physical equivalence
principle closing the amplitude \\
separate homogeneous conservation & effective-sector assumption pending
covariant completion \\
\(\Lambda_{\rm res}=0\) & global sector choice \\
spatial flatness & empirical/background sector choice \\
\end{longtable}

It is therefore not correct to say that mathematics alone derives every
premise. It is correct to say that one compact equivalence law replaces
the former free amplitude and that the resulting background history
contains no adjustable dark function.

\subsection{Stopping condition}\label{stopping-condition}

A foundational theory is not required to derive every axiom from
something deeper. It is adequate when:

\begin{enumerate}
\def\labelenumi{\arabic{enumi}.}
\tightlist
\item
  primitive quantities are operationally and mathematically typed;
\item
  independent representations are connected by explicit equivalence
  laws;
\item
  no hidden functions are introduced to fit individual anomalies;
\item
  the laws produce many linked predictions;
\item
  the theory states clear failure conditions;
\item
  further elaboration concerns new regimes, not repair of the background
  law.
\end{enumerate}

By this standard, the homogeneous theory has reached a defensible
stopping point:

\[
\boxed{
\left.\frac{k_B}{S_c}G^{\perp}_{NN}\right|_{N_c}=1
}
\]

is taken as a fundamental scale-capacity equivalence principle. The
background equations and amplitude follow. The perturbation lift and
global vacuum sector are separate layers, not unresolved coefficients in
the homogeneous model.

\section{Audit of the proposed final
closure}\label{audit-of-the-proposed-final-closure}

The attached AI closure contains a valuable synthesis but also two
invalid deductions. This section records the audit so that the final
theory does not inherit attractive category errors.

\subsection{Useful conclusions
retained}\label{useful-conclusions-retained}

The closure correctly emphasizes:

\begin{enumerate}
\def\labelenumi{\arabic{enumi}.}
\tightlist
\item
  vertical and horizontal modular sectors must remain distinct;
\item
  the BKM Hessian of a modular rescaling family is a capacity/variance;
\item
  a self-dual crossing is the natural place to evaluate the response;
\item
  the homogeneous theory should end with one structural law rather than
  an indefinitely expanding assumption ledger;
\item
  the perturbation sector remains open even after background closure.
\end{enumerate}

These points are incorporated into the present formulation.

\subsection{Rejected claim: conformal-weight
integrality}\label{rejected-claim-conformal-weight-integrality}

The proposed argument stated that

\[
\mathcal E[1]\oplus\mathcal E[-1]
\]

forces

\[
\varrho_\perp\in\mathbb Z,
\]

and that an observational existence ceiling then selects
\(\varrho_\perp=1\).

That does not follow. Conformal density bundles \(\mathcal E[w]\) exist
for real weights \(w\); the representation theory of the positive scale
group admits continuous real characters. There is no general integrality
theorem.

The valid statement is narrower:

\begin{quote}
The fundamental null-normal pair and the fundamental scale/inverse-scale
pair both carry characters \(e^{\pm x}\). Identifying those fundamental
representations gives \(\varrho_\perp=1\) as a physical representation
choice.
\end{quote}

The existence ceiling is a prediction and consistency bound. It cannot
be used to derive a universal constant from measured \(\Omega_m\)
without making the derivation data dependent and circular.

\subsection{Rejected claim: the normal plane is a two-dimensional
CFT}\label{rejected-claim-the-normal-plane-is-a-two-dimensional-cft}

The closure argued that the normal plane is two-dimensional, so Cardy
thermodynamics gives \(C=S\) and hence the capacity ratio equals one.

A two-dimensional Lorentzian normal vector space is not, by that fact
alone, a \(1+1\) dimensional conformal quantum field theory. Cardy's
formula requires a CFT, a Hamiltonian, central charge, boundary
conditions, and a thermodynamic regime. The dimension of the normal
plane does not supply these structures.

The valid replacement is:

\begin{quote}
Capacity equals entropy in specific Einstein-holographic spherical
settings and is conjectured or supported for certain causal horizons.
This motivates a universality principle for the self-dual Einstein wall;
it is not a consequence of dimensionality alone.
\end{quote}

\subsection{Why the revised closure is
stronger}\label{why-the-revised-closure-is-stronger}

The final law does not require an integrality theorem or an assumed 2D
CFT:

\[
\boxed{
\mathfrak R_c
=\frac{k_B}{S_c}G^{\perp}_{NN}(N_c)=1.
}
\]

It is coordinate invariant and combines the slope and capacity
normalization into one measurable quantity. Existing capacity--entropy
results supply precedent, while Hawking--Friedmann thermodynamics
supplies the dimensional conversion.

The consequences are then algebraic:

\[
\Omega_{X,c}=\frac12,
\]

\[
\rho_X(N_c)=\rho_{\rm ordinary}(N_c),
\]

and in general spatial dimension

\[
\Omega_{X,c}=\frac1{d-1}.
\]

This is an explicitly stated equivalence principle with falsifiable
consequences, not a theorem obtained by smuggling in observed
cosmological parameters.

\section{Status, stopping principle, and research
programme}\label{status-stopping-principle-and-research-programme}

\subsection{What is closed}\label{what-is-closed}

Conditional on the stated equivalence principles, the homogeneous theory
now fixes:

\begin{itemize}
\tightlist
\item
  the mathematical type of scale;
\item
  the local trace-free gravitational source equation;
\item
  the horizontal binary state geometry;
\item
  affine state--scale soldering;
\item
  the unit shape coefficient;
\item
  the amplitude at the self-dual crossing;
\item
  the entire background density and equation-of-state history;
\item
  the sign change of acceleration;
\item
  the future coasting branch;
\item
  the differential invariant and other observational tests.
\end{itemize}

There is no arbitrary function \(w(z)\), no free dark-energy particle
potential, and no fitted late-time amplitude.

\subsection{What is not closed}\label{what-is-not-closed}

The following remain separate research layers:

\begin{enumerate}
\def\labelenumi{\arabic{enumi}.}
\tightlist
\item
  \textbf{FLRW causal-wall construction.} Build the scale-indexed local
  algebras and states for a dynamical FLRW apparent horizon and derive
  the horizontal cocycle directly.
\item
  \textbf{Perturbation lift.} Embed the canonical Witten pair into a
  covariant scalar/vector/tensor perturbation system.
\item
  \textbf{Global vacuum sector.} Derive or select the residual
  top-form/tractor-norm sector, including whether
  \(\Lambda_{\rm res}=0\).
\item
  \textbf{Initial state and arrow.} Explain the low-Weyl, low-variance
  boundary condition and the selection of the expanding orientation.
\item
  \textbf{Full likelihood.} Implement the perturbation completion in a
  Boltzmann solver and test CMB, lensing, growth, BAO, and supernovae
  jointly.
\end{enumerate}

These do not reopen the homogeneous dark-history amplitude. They decide
whether the framework extends beyond its closed background sector.

\subsection{Research priorities}\label{research-priorities}

\begin{enumerate}
\def\labelenumi{\arabic{enumi}.}
\tightlist
\item
  Construct the self-dual FLRW causal diamond and calculate
  \(G^{\perp}_{NN}/(S/k_B)\) from modular data rather than postulating
  it.
\item
  Derive the spacetime perturbation operator from deformations of the
  causal-cut family and compare it with the Witten pair.
\item
  Test the invariant nonparametrically from reconstructed \(w(z)\).
\item
  Test the predicted coincidence of equality and \(w=-1\) crossing.
\item
  Perform the direct Boltzmann likelihood, with the neutrino residual
  treated as a secondary structured test.
\item
  Study the dimension formula as a possible consistency principle rather
  than a claim that spatial dimensionality has already been derived.
\end{enumerate}

\subsection{Falsifiability}\label{falsifiability}

The homogeneous framework is false if any of the following is
established:

\begin{itemize}
\tightlist
\item
  the redshift-dependent invariant is not constant or not four;
\item
  the self-dual crossing fraction is not one half;
\item
  equality and \(w=-1\) crossing are significantly separated;
\item
  the required crossing date conflicts with flatness and the ordinary
  budget;
\item
  the dark history has multiple maxima or repeated crossing episodes;
\item
  a derived perturbation completion necessarily introduces unstable or
  arbitrarily tunable modes;
\item
  the FLRW wall capacity ratio is calculated and is inconsistent with
  unity.
\end{itemize}

The scale-tractor reformulation and local vacuum-blind quotient can
survive failure of the late-time capacity model; the theory is
deliberately modular in its logical dependencies.

\section{Conclusion}\label{conclusion}

Causal Scale Dynamics begins from a standard theorem-level distinction:
causal order determines conformal geometry but not physical scale. The
scale is a section \(\sigma\in\Gamma(\mathcal E[1])\), and its two-jet
is packaged by the scale tractor. In this language, trace-free stress is
the obstruction to parallel scale transport, while the scalar curvature
and cosmological lift occupy a separate norm channel.

The quantum register supplies a second geometry. A causal region and
state carry modular structure; relative entropy supplies a BKM Hessian.
The corrected theory distinguishes vertical modular automorphism flow
from horizontal change of the state family. Homogeneous normal chirality
gives a binary exponential family with

\[
\eta=\tanh\theta,
\qquad
G^{\mathrm{BKM}}_{\theta\theta}=\operatorname{sech}^2\theta.
\]

Cocycle composition makes the horizontal coordinate affine in Weyl
scale. The fundamental representation sets

\[
\theta=N-N_c.
\]

The remaining amplitude is closed by the Scale--Capacity Equivalence
Principle,

\[
\boxed{
\left.\frac{k_B}{S_c}G^{\perp}_{NN}\right|_{N_c}=1.
}
\]

Relative entropy converts this dimensionless response into modular
free-energy curvature, and horizon thermodynamics supplies the units:

\[
\rho_X=\frac{k_BT_c}{2V_c}G^{\perp}_{NN},
\]

\[
T_cS_c=E_{\rm MS,c}=\rho_{\rm crit,c}V_c.
\]

Therefore, in \(3+1\) dimensions,

\[
\boxed{
\rho_X(N)=\frac12\rho_{\rm crit,c}\operatorname{sech}^2(N-N_c).
}
\]

At self-duality,

\[
\boxed{
\rho_X=\rho_{\rm ordinary},
\qquad
w_X=-1,
\qquad
\Omega_X=\frac12.
}
\]

Separate conservation gives

\[
\boxed{
w_X=-1+\frac23\tanh(N-N_c),
}
\]

and the structural test

\[
\boxed{
9(1+w_X)^2+6w_X'=4.
}
\]

The observed acceleration transition follows from the changing balance
between the ordinary sector and this one self-dual response. It is not
generated by an independently postulated repulsive force. The same law
predicts a future exit and asymptotic coasting when the global residual
floor is exactly zero.

The theory has therefore reached a legitimate homogeneous stopping
point: one clearly stated equivalence principle replaces the former
amplitude postulate and produces a closed, economical background
history. The next problems are not further adjustment of that history.
They are the construction of the dynamical FLRW wall, the covariant
perturbation lift, and the global vacuum sector.

The central retyping is:

\[
\boxed{
\text{late cosmic response is the modular free-energy curvature of causal-state change under scale.}
}
\]

This statement is precise enough to inspire independent mathematical
construction and sufficiently constrained to be rejected by observation.

\appendix

\section{Symbol dictionary}\label{symbol-dictionary}

\begin{longtable}[]{@{}
  >{\raggedright\arraybackslash}p{(\columnwidth - 4\tabcolsep) * \real{0.3333}}
  >{\raggedright\arraybackslash}p{(\columnwidth - 4\tabcolsep) * \real{0.3333}}
  >{\raggedright\arraybackslash}p{(\columnwidth - 4\tabcolsep) * \real{0.3333}}@{}}
\toprule\noalign{}
\begin{minipage}[b]{\linewidth}\raggedright
symbol
\end{minipage} & \begin{minipage}[b]{\linewidth}\raggedright
definition
\end{minipage} & \begin{minipage}[b]{\linewidth}\raggedright
mathematical type / status
\end{minipage} \\
\midrule\noalign{}
\endhead
\bottomrule\noalign{}
\endlastfoot
\((M,[g])\) & conformal spacetime & standard Lorentzian conformal
geometry \\
\(\boldsymbol g\) & conformal metric & section of the weighted metric
bundle \\
\(\sigma\in\Gamma(\mathcal E[1])\) & positive scale & density-bundle
section \\
\(g_{\rm phys}=\sigma^{-2}\boldsymbol g\) & physical metric &
conformally invariant combination \\
\(I_A=\tfrac14D_A\sigma\) & scale tractor & packages the scale
two-jet \\
\(P_{ab}\) & Schouten tensor & standard conformal curvature tensor \\
\(N=-\ln(\sigma/\sigma_c)\) & Weyl e-fold coordinate & additive scale
displacement \\
\(s\) & vertical modular automorphism parameter & state-preserving
modular flow \\
\(\eta_A\) & vertical horizon rapidity potential & geometric definition;
modular interpretation conditional \\
\(\theta\) & horizontal state coordinate & noncentral relative modular
polarization \\
\(Q=P_+-P_-\) & normal chirality & binary quotient, \(Q^2=1\),
\(JQJ=-Q\) \\
\(\Psi(\theta)=\ln(2\cosh\theta)\) & log-partition potential & binary
exponential family \\
\(\eta=\langle Q\rangle\) & mixture coordinate & \(\eta=\tanh\theta\) \\
\(G^{\rm BKM}\) & BKM metric & Hessian of Umegaki relative entropy \\
\(G^{\perp}_{NN}\) & pullback BKM norm & extensive horizontal capacity
per squared e-fold \\
\(S_c\) & self-dual wall entropy & Bekenstein--Hawking entropy in the
closure law \\
\(T_c\) & horizontal modular temperature & canonical \(2\pi\)
normal-boost normalization \\
\(V_c\) & causal-wall volume & areal volume of the crossing apparent
horizon \\
\(\mathfrak R_c\) & scale-capacity number &
\((k_B/S_c)G^{\perp}_{NN}(N_c)\); fundamental value \(1\) \\
\(\varrho_\perp\) & horizontal soldering slope & \(d\theta/dN\);
fundamental representation value \(1\) \\
\(\rho_X,p_X,w_X\) & effective homogeneous response variables & derived
from the scale-capacity source and conservation \\
\(N_c,z_c\) & self-dual crossing location & intrinsic event, cosmic date
fixed by the solution \\
\(\mu_A\) & horizon modular--Weyl index & \((1-q)/2=-I^2/H^2\); running
state variable \\
\(\Lambda_g,\Lambda_{\rm res}\) & scalar lift/global residual & global
sector, not local trace-free source \\
\(\mathcal A,\mathcal A^\dagger\) & Witten/Darboux first-order operators
& generate the horizontal pair \\
\(\mathcal H_\pm\) & horizontal partner operators & one reflectionless
Pöschl--Teller partner and one free partner \\
\end{longtable}

\section{Core derivations}\label{core-derivations}

\subsection{Binary information
geometry}\label{binary-information-geometry-1}

For

\[
\omega_\theta=\frac{e^{\theta Q}}{2\cosh\theta},
\qquad Q^2=1,
\]

the partition potential is

\[
\Psi=\ln\operatorname{tr}e^{\theta Q}=\ln(2\cosh\theta).
\]

Therefore

\[
\Psi'=\tanh\theta,
\qquad
\Psi''=\operatorname{sech}^2\theta,
\]

and

\[
\langle Q\rangle^2+\operatorname{Var}(Q)=1.
\]

\subsection{Cocycle soldering}\label{cocycle-soldering}

Assume the reduced noncentral Connes cocycle depends on two scales only
through \(r=\sigma_2/\sigma_1\) and is measurable. Cocycle composition
gives

\[
\theta(r_1r_2)=\theta(r_1)+\theta(r_2).
\]

The measurable solutions are

\[
\theta(r)=-\varrho_\perp\ln r.
\]

Since \(N=-\ln(\sigma/\sigma_c)\),

\[
\theta=\varrho_\perp(N-N_c).
\]

The equation fixes the form but not the value of \(\varrho_\perp\).

\subsection{Free-energy Hessian}\label{free-energy-hessian}

For a KMS reference state with physical modular Hamiltonian
\(\mathcal H_c=k_BT_cK_c\),

\[
F_c(\rho)-F_c(\omega_c)=k_BT_cS(\rho\|\omega_c).
\]

At coincidence,

\[
S(\omega_{c+\delta N}\|\omega_c)
=\frac12G^{\perp}_{NN}(N_c)\delta N^2+O(\delta N^3).
\]

Therefore the quadratic free-energy curvature per causal-wall volume is

\[
\rho_{X,c}=\frac{k_BT_c}{2V_c}G^{\perp}_{NN}(N_c).
\]

\subsection{Hawking--Friedmann
conversion}\label{hawkingfriedmann-conversion-1}

In \(3+1\) dimensional flat FLRW,

\[
R_c=\frac{c}{H_c},
\quad
\frac{S_c}{k_B}=\frac{\pi R_c^2c^3}{G\hbar},
\quad
k_BT_c=\frac{\hbar c}{2\pi R_c}.
\]

Thus

\[
k_BT_c\frac{S_c}{k_B}=\frac{c^4R_c}{2G}=E_{\rm MS,c}.
\]

Since

\[
V_c=\frac{4\pi R_c^3}{3},
\]

\[
\frac{E_{\rm MS,c}}{V_c}=\frac{3c^4}{8\pi GR_c^2}=\frac{3c^2H_c^2}{8\pi G}=\rho_{\rm crit,c}.
\]

\subsection{Scale-capacity amplitude}\label{scale-capacity-amplitude}

The equivalence law is

\[
G^{\perp}_{NN}(N_c)=\mathfrak R_c\frac{S_c}{k_B}.
\]

Therefore

\[
\rho_{X,c}=\frac{\mathfrak R_c}{2}\rho_{\rm crit,c}.
\]

For \(\mathfrak R_c=1\), flatness gives

\[
\rho_{X,c}=\rho_{\rm ordinary,c}=\frac12\rho_{\rm crit,c}.
\]

\subsection{General dimension}\label{general-dimension}

For \(d\) spatial dimensions,

\[
A=\Omega_{d-1}R^{d-1},
\qquad
V=\frac{\Omega_{d-1}}{d}R^d,
\]

and Einstein--Friedmann marginality gives

\[
\frac{k_BT(S/k_B)}{V}=\frac{2}{d-1}\rho_{\rm crit}.
\]

Thus

\[
\Omega_{X,c}=\frac{\mathfrak R_c}{d-1},
\]

and

\[
\frac{\rho_{X,c}}{\rho_{\rm ordinary,c}}=
\frac{\mathfrak R_c}{d-1-\mathfrak R_c}.
\]

\subsection{Shape invariant}\label{shape-invariant}

With

\[
\rho_X=\rho_*\operatorname{sech}^2(N-N_c),
\]

\[
\Delta_X:=-\frac{\mathrm d\ln\rho_X}{\mathrm dN}=2\tanh(N-N_c).
\]

Then

\[
\Delta_X'=2\operatorname{sech}^2(N-N_c)=2-\frac12\Delta_X^2,
\]

so

\[
\Delta_X^2+2\Delta_X'=4.
\]

Using \(\Delta_X=3(1+w_X)\) gives

\[
9(1+w_X)^2+6w_X'=4.
\]

\subsection{Witten pair}\label{witten-pair}

Let \(\eta=\tanh\theta\). Then

\[
\eta^2+\eta'=1.
\]

With

\[
\mathcal A=\partial_\theta+\eta,
\qquad
\mathcal A^\dagger=-\partial_\theta+\eta,
\]

\[
\mathcal A^\dagger\mathcal A=-\partial_\theta^2+1-2\operatorname{sech}^2\theta,
\]

\[
\mathcal A\mathcal A^\dagger=-\partial_\theta^2+1.
\]

The normalized zero mode is
\(\psi_0=2^{-1/2}\operatorname{sech}\theta\).

\section{Receipt, data, and audit
ledger}\label{receipt-data-and-audit-ledger}

The accompanying script \texttt{receipts/receipts\_v7.py} verifies the
exact identities and benchmark values used in this note. The
machine-readable output is \texttt{receipts/receipts\_v7.json}; the
expected field is

\begin{Shaded}
\begin{Highlighting}[]
\ErrorTok{"all\_exact\_residuals\_zero":} \ErrorTok{true}
\end{Highlighting}
\end{Shaded}

The receipts verify:

\begin{itemize}
\tightlist
\item
  binary moments and the BKM metric;
\item
  affine-shape identities;
\item
  Fisher length \(\pi\);
\item
  the Witten/Darboux factorization and zero mode;
\item
  Hawking--Friedmann conversion;
\item
  the dimension-dependent crossing fraction;
\item
  the closed benchmark and radiation correction;
\item
  present cosmography and the two acceleration transitions.
\end{itemize}

The package includes the earlier background-fit, direct-response, and
economy-audit results used for the observational status section. These
analyses use different baselines and are not combined as independent
measurements.

The attached AI referee notes were used as adversarial inputs. Their
useful results and rejected claims are recorded in the body. In
particular:

\begin{itemize}
\tightlist
\item
  the cocycle argument for affine soldering was retained;
\item
  the vertical/horizontal distinction was retained;
\item
  the warnings about the sigma-model completion and response
  normalization were retained;
\item
  conformal-weight integrality and the identification of a
  two-dimensional normal plane with a two-dimensional CFT were rejected.
\end{itemize}

\section{Epistemic status ledger}\label{epistemic-status-ledger}

\begin{longtable}[]{@{}
  >{\raggedright\arraybackslash}p{(\columnwidth - 2\tabcolsep) * \real{0.5000}}
  >{\raggedright\arraybackslash}p{(\columnwidth - 2\tabcolsep) * \real{0.5000}}@{}}
\toprule\noalign{}
\begin{minipage}[b]{\linewidth}\raggedright
statement
\end{minipage} & \begin{minipage}[b]{\linewidth}\raggedright
status
\end{minipage} \\
\midrule\noalign{}
\endhead
\bottomrule\noalign{}
\endlastfoot
causal order fixes conformal geometry under standard hypotheses &
standard theorem \\
scale as \(\Gamma(\mathcal E[1])\) and scale tractor & standard
conformal geometry \\
Einstein metric iff parallel scale tractor & standard theorem \\
trace-free tractor source equation & exact reformulation of GR \\
modular/gravitational boost-charge identity & physical unity principle,
supported in controlled settings \\
fundamental normal chirality quotient & structural identification \\
affine cocycle law & conditional theorem from ratio dependence and
measurability \\
\(\varrho_\perp=1\) & fundamental representation choice \\
binary BKM metric \(\operatorname{sech}^2\theta\) & exact after
chirality reduction \\
scale--capacity number \(\mathfrak R_c=1\) & fundamental equivalence
principle \\
free-energy source law & constitutive definition motivated by relative
entropy \\
Hawking--Friedmann conversion & exact horizon/Friedmann identity in the
stated regime \\
closed homogeneous pulse & deduction from the preceding laws \\
Witten pair & exact internal horizontal construction \\
covariant perturbation theory & open \\
zero residual floor & global sector choice \\
sequestering completion & published candidate, not derived from the
scale-capacity law \\
\end{longtable}

\section{References}\label{references}

\begin{enumerate}
\def\labelenumi{\arabic{enumi}.}
\tightlist
\item
  S. W. Hawking, A. R. King, and P. J. McCarthy, ``A new topology for
  curved space-time which incorporates the causal, differential, and
  conformal structures,'' \emph{J. Math. Phys.} \textbf{17} (1976) 174.
\item
  D. B. Malament, ``The class of continuous timelike curves determines
  the topology of spacetime,'' \emph{J. Math. Phys.} \textbf{18} (1977)
  1399.
\item
  T. N. Bailey, M. G. Eastwood, and A. R. Gover, ``Thomas's structure
  bundle for conformal, projective and related structures,'' \emph{Rocky
  Mountain J. Math.} \textbf{24} (1994) 1191.
\item
  S. Curry and A. R. Gover, ``An introduction to conformal geometry and
  tractor calculus, with a view to applications in general relativity,''
  \href{https://arxiv.org/abs/1412.7559}{arXiv:1412.7559}.
\item
  A. R. Gover, ``Almost Einstein and Poincaré--Einstein manifolds in
  Riemannian signature,'' \emph{J. Geom. Phys.} \textbf{60} (2010) 182.
\item
  J. J. Bisognano and E. H. Wichmann, ``On the duality condition for
  quantum fields,'' \emph{J. Math. Phys.} \textbf{17} (1976) 303.
\item
  H. Casini, M. Huerta, and R. C. Myers, ``Towards a derivation of
  holographic entanglement entropy,''
  \href{https://arxiv.org/abs/1102.0440}{arXiv:1102.0440}.
\item
  R. M. Wald, ``Black hole entropy is the Noether charge,''
  \href{https://arxiv.org/abs/gr-qc/9307038}{arXiv:gr-qc/9307038}.
\item
  T. Jacobson, ``Thermodynamics of spacetime: The Einstein equation of
  state,''
  \href{https://arxiv.org/abs/gr-qc/9504004}{arXiv:gr-qc/9504004}.
\item
  T. Jacobson, ``Entanglement equilibrium and the Einstein equation,''
  \href{https://arxiv.org/abs/1505.04753}{arXiv:1505.04753}.
\item
  D. L. Jafferis, A. Lewkowycz, J. Maldacena, and S. J. Suh, ``Relative
  entropy equals bulk relative entropy,''
  \href{https://arxiv.org/abs/1512.06431}{arXiv:1512.06431}.
\item
  N. Lashkari and M. Van Raamsdonk, ``Canonical energy is quantum Fisher
  information,''
  \href{https://arxiv.org/abs/1508.00897}{arXiv:1508.00897}.
\item
  B. Czech, L. Lamprou, S. McCandlish, and J. Sully, ``Modular Berry
  Connection,''
  \href{https://arxiv.org/abs/1712.07123}{arXiv:1712.07123}.
\item
  B. Czech et al., ``Changing states in holography: From modular Berry
  curvature to the bulk symplectic form,''
  \href{https://arxiv.org/abs/2305.16384}{arXiv:2305.16384}.
\item
  D. Petz, ``Monotone metrics on matrix spaces,'' \emph{Linear Algebra
  Appl.} \textbf{244} (1996) 81.
\item
  N. N. Čencov, \emph{Statistical Decision Rules and Optimal Inference},
  AMS (1982).
\item
  S.-I. Amari and H. Nagaoka, \emph{Methods of Information Geometry},
  AMS/Oxford (2000).
\item
  M. R. Grasselli and R. F. Streater, ``On the uniqueness of the
  Chentsov metric in quantum information geometry,''
  \href{https://arxiv.org/abs/math-ph/0006030}{arXiv:math-ph/0006030}.
\item
  R. Chatterjee, ``Modular Self-Duality, Symmetrized Relative Entropy,
  and Bogoliubov--Kubo--Mori Susceptibility in Quantum Field Theory,''
  \href{https://arxiv.org/abs/2605.19106}{arXiv:2605.19106}.
\item
  K. Jensen, J. Sorce, and A. Speranza, ``Generalized entropy for
  general subregions in quantum gravity,''
  \href{https://arxiv.org/abs/2306.01837}{arXiv:2306.01837}.
\item
  T. Faulkner and A. J. Speranza, ``Gravitational algebras and the
  generalized second law,''
  \href{https://arxiv.org/abs/2405.00847}{arXiv:2405.00847}.
\item
  V. Chandrasekaran and É. É. Flanagan, ``Subregion algebras in
  classical and quantum gravity,''
  \href{https://arxiv.org/abs/2601.07915}{arXiv:2601.07915}.
\item
  S. A. Hayward, ``Unified first law of black-hole dynamics and
  relativistic thermodynamics,''
  \href{https://arxiv.org/abs/gr-qc/9710089}{arXiv:gr-qc/9710089}.
\item
  R.-G. Cai and S. P. Kim, ``First law of thermodynamics and Friedmann
  equations of the Friedmann--Robertson--Walker universe,''
  \href{https://arxiv.org/abs/hep-th/0501055}{arXiv:hep-th/0501055}.
\item
  M. Akbar and R.-G. Cai, ``Thermodynamic behavior of Friedmann equation
  at apparent horizon of FRW universe,''
  \href{https://arxiv.org/abs/hep-th/0609128}{arXiv:hep-th/0609128}.
\item
  J. de Boer, J. Järvelä, and E. Keski-Vakkuri, ``Aspects of capacity of
  entanglement,''
  \href{https://arxiv.org/abs/1807.07357}{arXiv:1807.07357}.
\item
  T. Banks and K. M. Zurek, ``Conformal description of near-horizon
  vacuum states,'' \emph{Phys. Rev.~D} \textbf{104}, 126026 (2021),
  \href{https://arxiv.org/abs/2108.04806}{arXiv:2108.04806}.
\item
  T. Banks and P. Draper, ``Generalized entanglement capacity of de
  Sitter space,''
  \href{https://arxiv.org/abs/2404.13684}{arXiv:2404.13684}.
\item
  D. Kastor, S. Ray, and J. Traschen, ``Enthalpy and the mechanics of
  AdS black holes,''
  \href{https://arxiv.org/abs/0904.2765}{arXiv:0904.2765}.
\item
  N. Kaloper, A. Padilla, D. Stefanyszyn, and G. Zahariade, ``A
  manifestly local theory of vacuum energy sequestering,''
  \href{https://arxiv.org/abs/1505.01492}{arXiv:1505.01492}.
\item
  N. Kaloper and A. Padilla, ``Vacuum energy sequestering and graviton
  loops,'' \href{https://arxiv.org/abs/1606.04958}{arXiv:1606.04958}.
\item
  A. Vikman, ``Can dark energy evolve to the phantom?''
  \href{https://arxiv.org/abs/astro-ph/0407107}{arXiv:astro-ph/0407107}.
\item
  J. Lekner, ``Reflectionless eigenstates of the
  \(\operatorname{sech}^2\) potential,'' \emph{Am. J. Phys.} \textbf{75}
  (2007) 1151.
\item
  DESI Collaboration, ``DESI DR2 results: measurements of baryon
  acoustic oscillations and cosmological constraints,''
  \href{https://arxiv.org/abs/2503.14738}{arXiv:2503.14738}.
\item
  DESI Collaboration, ``DESI DR2 results: neutrino mass constraints,''
  \href{https://arxiv.org/abs/2503.14744}{arXiv:2503.14744}.
\item
  DESI Collaboration, ``DESI DR2 Results IV: Ly\(\alpha\)-forest
  full-shape measurements and cosmological constraints,''
  \href{https://arxiv.org/abs/2607.27410}{arXiv:2607.27410}.
\item
  D. Brout et al., ``The Pantheon+ analysis: cosmological constraints,''
  \href{https://arxiv.org/abs/2202.04077}{arXiv:2202.04077}.
\item
  X. Li and A. Shafieloo, ``A simple phenomenological emergent dark
  energy model,''
  \href{https://arxiv.org/abs/1906.08275}{arXiv:1906.08275}.
\item
  L. Parker and A. Raval, ``Vacuum-driven metamorphosis,''
  \href{https://arxiv.org/abs/gr-qc/0312108}{arXiv:gr-qc/0312108}.
\item
  J. Solà Peracaula et al., ``Running vacuum in the Universe,''
  \href{https://arxiv.org/abs/2203.13757}{arXiv:2203.13757}.
\end{enumerate}

\end{document}
